\documentclass[a4paper,10pt]{article}
\usepackage[utf8]{inputenc}
\usepackage[T1]{fontenc}
\usepackage[italian]{babel}
\usepackage{lmodern}
\usepackage[margin=1.7cm]{geometry}
\usepackage{bpninja}
\usepackage{booktabs,tabularx,longtable,array,enumitem}
\usepackage[hidelinks]{hyperref}
\usepackage{xcolor}
\setlist{itemsep=1pt,topsep=3pt}
\setlength{\parindent}{0pt}
\setlength{\parskip}{4pt}
\sloppy\emergencystretch=3em
\newcommand{\nodo}[1]{\textsf{\textbf{#1}}}
\newcommand{\var}[1]{\texttt{#1}}
\newcommand{\tasto}[1]{\fbox{\footnotesize\textsf{#1}}}
\newcommand{\ninja}[1]{\par\smallskip\noindent\colorbox{black!6}{\parbox{\dimexpr\linewidth-2\fboxsep}{\small\textbf{Mossa da ninja.} #1}}\par\smallskip}
\newcommand{\attenzione}[1]{\par\smallskip\noindent\colorbox{orange!12}{\parbox{\dimexpr\linewidth-2\fboxsep}{\small\textbf{Attenzione.} #1}}\par\smallskip}
\newcolumntype{L}[1]{>{\raggedright\arraybackslash}p{#1}}

% tabella degli oggetti di un livello
\newenvironment{oggetti}{\par\footnotesize\setlength{\tabcolsep}{3pt}\renewcommand{\arraystretch}{1.15}%
\begin{longtable}{@{}r L{2.35cm} L{2.9cm} L{2.25cm} L{2.4cm} L{5.85cm}@{}}
\toprule
\# & Cosa & Nome nell'Outliner & Posizione (X; Y; Z) & Misure / rotazione & Impostazioni \\
\midrule\endhead
\bottomrule\endfoot}{\end{longtable}\normalsize}
\newcommand{\nota}[1]{\multicolumn{6}{@{}L{17.2cm}@{}}{\color{black!60}\itshape #1}\\}

% tabella delle variabili di un Blueprint
\newenvironment{variabili}{\par\footnotesize\setlength{\tabcolsep}{3pt}%
\begin{longtable}{@{}L{3.3cm} L{3cm} L{2.3cm} L{8.2cm}@{}}
\toprule Variabile & Tipo & Default & A cosa serve \\ \midrule\endhead
\bottomrule\endfoot}{\end{longtable}\normalsize}

\title{\textbf{Rifare i livelli in Blueprint}\\[4pt]\large Guida completa per ricostruire \texttt{L\_GD13\_Livello1}, \texttt{2} e \texttt{3} senza una riga di C++, e il menu che li sceglie}
\author{Progetto \texttt{AIV\_UE5Intro}}
\date{}

\begin{document}
\maketitle
\tableofcontents
\newpage

% =====================================================================
\section{Il menu e le due versioni}

Il gioco ora parte da un livello ``padre'', \var{L\_GD13\_Menu}. È una piccola scena al tramonto vista da una telecamera fissa; sopra, due pulsanti:

\begin{itemize}
  \item \textbf{C++}: apre \var{L\_GD13\_Livello1}, i livelli costruiti con le classi C++ del progetto (\texttt{Source/AIV\_UE5Intro/Platformer});
  \item \textbf{Blueprints}: apre \var{L\_GD13\_Livello1\_BP}, la copia esatta che costruirete voi con questa guida.
\end{itemize}

Si sceglie con il mouse, oppure con \tasto{1}/\tasto{C} e \tasto{2}/\tasto{B}, le frecce e \tasto{Invio}, o con il pad. Finché \var{L\_GD13\_Livello1\_BP} non esiste, il pulsante Blueprints è arancione e dice ``non ancora creato'': cliccandolo compare un avviso invece di aprire una mappa vuota. Nei livelli C++ il tasto \tasto{M} riporta al menu, e il traguardo del livello 3 riporta al menu invece che al livello 1.

\subsection{I nomi dei livelli Blueprint}

\begin{center}
\begin{tabular}{lll}
\toprule
Versione C++ & Versione Blueprint & Si va avanti con \\
\midrule
\var{L\_GD13\_Livello1} & \var{L\_GD13\_Livello1\_BP} & traguardo $\rightarrow$ \var{L\_GD13\_Livello2\_BP} \\
\var{L\_GD13\_Livello2} & \var{L\_GD13\_Livello2\_BP} & traguardo $\rightarrow$ \var{L\_GD13\_Livello3\_BP} \\
\var{L\_GD13\_Livello3} & \var{L\_GD13\_Livello3\_BP} & traguardo $\rightarrow$ \var{L\_GD13\_Menu} \\
\bottomrule
\end{tabular}
\end{center}

Il suffisso è \var{\_BP} (maiuscolo, con il trattino basso) e non ``bp'' attaccato: nel Content Browser ogni livello Blueprint compare subito sotto il suo gemello C++, ed è la stessa convenzione dei prefissi che usa Unreal (\var{BP\_}, \var{MI\_}, \var{SM\_}). Tutti e tre vanno salvati in \texttt{Content/Maps}, cioè \texttt{/Game/Maps}: il menu cerca proprio \texttt{/Game/Maps/L\_GD13\_Livello1\_BP}. Se un giorno preferite un altro nome, cambiatelo nel menu (variabile \var{BlueprintFirstLevel} dell'HUD del menu, \texttt{PlatformerMenu.h}).

\subsection{Creare il menu}

Il menu lo crea lo stesso script dei livelli C++: in Unreal, \nodo{Tools} $\rightarrow$ \nodo{Execute Python Script\dots} $\rightarrow$ \texttt{Content/Python/gd13\_platformer.py} (che ora costruisce il menu e i tre livelli C++). Se volete solo il menu: nel \emph{Output Log}, riga di comando in modalità \emph{Python},
\begin{center}\texttt{import gd13\_platformer as g; g.M.update(g.mat.tutti()); g.menu()}\end{center}
Il livello usa il GameMode C++ \var{PlatformerMenuGameMode}: nessun personaggio, mouse visibile, l'HUD disegna i pulsanti. Perché il gioco parta dal menu anche in \emph{Standalone} o nel gioco impacchettato, in \nodo{Project Settings} $\rightarrow$ \nodo{Maps \& Modes} impostate \emph{Game Default Map} = \var{L\_GD13\_Menu} (l'\emph{Editor Startup Map} potete lasciarla com'è).

\subsection{Tornare al menu e il game over nella versione Blueprint}

Nei livelli \var{\_BP} non c'è niente di C++, quindi due piccole cose le fate voi:

\begin{itemize}
  \item \textbf{Il tasto M.} In \var{BP\_Eroe}: un evento di tastiera \nodo{M} (o una Input Action \var{IA\_Menu} nel vostro \var{IMC\_Ninja}) $\rightarrow$ \nodo{Open Level (by Name)} \var{L\_GD13\_Menu}.
  \item \textbf{Il primo livello del game over.} In \var{GI\_Platform} la variabile \var{PrimoLivello} vale \var{L\_GD13\_Livello1\_BP}: a zero vite si ricomincia dal primo livello \emph{della vostra} versione, non da quello C++ (guida Ninja, capitolo 2).
\end{itemize}

\begin{bpgraph}
\bpnode{k}{(0,0)}{bpev}{M (Keyboard Event)}{-}{exec/Pressed,exec/Released,struct/Key}
\bpnode{c}{\dx{k}{0.7}{0}}{bpfn}{Open Level (by Name)}{exec/,name/Level Name: L\_GD13\_Menu,str/Options}{exec/}
\bpexec{k-out1}{c-in1}
\end{bpgraph}

Il menu passa per \var{UPlatformerRunSubsystem::StartMode}: se volete sapere da Blueprint in quale versione si sta giocando, c'è il nodo \nodo{Is Blueprint Mode}.

% =====================================================================
\section{Prima di cominciare}

\subsection{Quello che serve sapere dell'editor}

\begin{itemize}
  \item \textbf{Muoversi nel viewport}: tasto destro tenuto premuto + \tasto{W}\tasto{A}\tasto{S}\tasto{D} (\tasto{Q}/\tasto{E} giù e su), rotellina per la velocità. \tasto{F} inquadra l'actor selezionato. \tasto{G} nasconde le icone (vista ``di gioco'').
  \item \textbf{Details}: il pannello a destra. Le caselle \emph{Location}, \emph{Rotation}, \emph{Scale} si scrivono a mano: è così che si mettono i numeri di questa guida. Cliccando sul lucchetto accanto a \emph{Scale} si sbloccano i tre valori.
  \item \textbf{Outliner}: l'elenco degli actor del livello. \tasto{F2} rinomina; tasto destro $\rightarrow$ \emph{Move To} $\rightarrow$ \emph{Create New Folder} per le cartelle. Usate gli stessi nomi delle tabelle: quando confronterete la vostra versione con quella C++ ritroverete ogni pezzo.
  \item \textbf{Duplicare}: \tasto{Ctrl}+\tasto{D}, oppure \tasto{Alt} mentre si trascina una freccia del gizmo. Dopo, correggete i numeri nei Details.
  \item \textbf{Provare}: \nodo{Play} (\tasto{Alt}+\tasto{P}) parte dal Player Start; tasto destro nel viewport $\rightarrow$ \emph{Play From Here} parte dal punto inquadrato: comodissimo per provare un tratto senza rifare tutto il livello.
\end{itemize}

\subsection{Come leggere le coordinate}

Tutto è in centimetri. Il percorso va sempre verso \textbf{+X}; +Y è a destra guardando il percorso, +Z è in alto. Nelle tabelle:
\begin{itemize}
  \item \textbf{Posizione} è il \emph{centro} dell'oggetto per i blocchi; per i Blueprint è l'origine dell'actor, descritta nel catalogo (per esempio, la torretta ha l'origine \emph{a terra}, la piattaforma mobile al \emph{centro}).
  \item \textbf{Misure} di un blocco sono le dimensioni intere in cm: lunghezza X, larghezza Y, altezza Z. Con il \var{Cube} dell'engine (che è largo 100\,cm) la \emph{Scale} è la misura divisa per 100: un blocco $400 \times 400 \times 40$ ha Scale $(4;\ 4;\ 0{,}4)$.
  \item Una \textbf{pedana} è un blocco su cui si cammina: la tabella dice anche la quota del piano (``si cammina a Z = \dots''). La Z della posizione è già il centro (piano meno metà spessore).
  \item \textbf{yaw} è la rotazione intorno a Z (Details $\rightarrow$ Rotation, terzo valore).
\end{itemize}

\subsection{I blocchi: Cube dell'engine o BlockOut Tools}

I livelli C++ usano il \var{Blockout\_Box} del plugin BlockOut Tools (già attivo nel progetto). Per voi è più semplice e identico nelle misure usare il \textbf{Cube dell'engine} (\nodo{Place Actors} $\rightarrow$ \nodo{Shapes} $\rightarrow$ \nodo{Cube}): la sua origine è il centro, quindi la posizione delle tabelle va scritta così com'è. Il colore lo dà un materiale già pronto in \texttt{/Game/AIV/Materiali}:

\begin{center}
\begin{tabular}{lll}
\toprule
Colore nei livelli C++ & Significato & Materiale da usare \\
\midrule
grigio chiaro & percorso & \var{MI\_AIV\_Percorso} \\
grigio scuro & percorso (alternato) & \var{MI\_AIV\_PercorsoScuro} \\
arancio & si salta / si interagisce & \var{MI\_AIV\_Salto} \\
blu & muro del wall run (tag \var{WallRun}) & \var{MI\_AIV\_Muro} \\
rosso & pericolo (colonne delle torrette) & \var{MI\_AIV\_Pericolo} \\
verde & traguardo & \var{MI\_AIV\_Traguardo} \\
marrone & ripari & \var{MI\_AIV\_LegnoScuro} \\
\bottomrule
\end{tabular}
\end{center}

Il materiale si mette trascinandolo dal Content Browser sul cubo, o in Details $\rightarrow$ \emph{Materials} $\rightarrow$ \emph{Element 0}. I materiali \var{MI\_AIV\_*} sono ``triplanari'': la griglia ha la stessa misura su ogni faccia comunque scaliate il cubo. Se preferite il Blockout\_Box: impostate \emph{Box Size} con le misure della tabella e controllate con lo strumento \emph{Measure} (tasto centrale del mouse trascinato nel viewport ortografico) che il centro sia dove deve: il suo pivot non è al centro.

\attenzione{I muri blu del wall run devono avere il \textbf{tag} \var{WallRun}: Details $\rightarrow$ cercate ``Tags'' $\rightarrow$ \emph{Actor} $\rightarrow$ \emph{Tags} $\rightarrow$ \nodo{+} $\rightarrow$ \var{WallRun}. Senza tag il personaggio non ci corre sopra (guida Ninja, capitolo 3.4).}

\subsection{Gli altri asset}

Sono tutti già nel progetto: i materiali in \texttt{/Game/AIV/Materiali} (\var{MI\_AIV\_Erba}, \var{\_Acqua}, \var{\_Lava}, \var{\_Moneta}, \var{\_Abilita}\dots), i suoni in \texttt{/Game/Audio} (\var{SfxPlatformMoving\_Cue}, \var{Rock\_Cue}), le piante in \texttt{/Game/Meshes} (\var{SM\_Chestnut\_4}, \var{SM\_Groundplant\_12}), le forme in \texttt{/Engine/BasicShapes} (per vederle nel Content Browser: \emph{Settings} $\rightarrow$ \emph{Show Engine Content}).

\subsection{I Blueprint del personaggio (dalla guida Ninja)}

I livelli Blueprint hanno bisogno del vostro personaggio e della vostra ``memoria'' del gioco. Sono tutti spiegati nodo per nodo in \texttt{Docs/Blueprint\_Ninja}:

\begin{center}
\begin{tabularx}{\linewidth}{lX}
\toprule
Asset & Capitolo della guida Ninja \\
\midrule
\var{GI\_Platform} (vite, monete, abilità, salvataggio) & 2 \\
\var{BP\_Eroe} con \var{AC\_Movimento} (coyote, doppio salto, wall run, jetpack) & 3 \\
\var{AC\_Torcia} (prima persona con la rotella, batteria) & 4 \\
\var{AC\_PostoSicuro} (tre vite, ritorno al posto sicuro) & 5 \\
\var{AC\_Casse} e \var{BP\_Cassa} (generatore di casse) & 6 \\
\var{WBP\_HUD} & 12 \\
\var{GM\_Platform}: Default Pawn Class = \var{BP\_Eroe}, HUD/Widget & 1 \\
\var{BPI\_Interagibile}, \var{BPI\_Attivabile} (interfacce) & 9 \\
\bottomrule
\end{tabularx}
\end{center}

Il consiglio è di costruire il personaggio \emph{prima} dei livelli, almeno movimento e vite: un livello senza un personaggio che salta come quello C++ non si può collaudare (le distanze sono tarate su salto 650, doppio salto 600, wall run 650).

\subsection{Dove mettere i file}

\begin{center}
\begin{tabular}{ll}
\toprule
Cartella & Cosa \\
\midrule
\texttt{Content/Maps} & \var{L\_GD13\_Livello1\_BP}, \var{\_2\_BP}, \var{\_3\_BP} \\
\texttt{Content/Blueprints/Livelli/Zone} & \var{BP\_Zona} e figli \\
\texttt{Content/Blueprints/Livelli/Pickup} & \var{BP\_Pickup} e figli \\
\texttt{Content/Blueprints/Livelli/Congegni} & interruttori, attivabili, piattaforme, trampolino \\
\texttt{Content/Blueprints/Livelli/Ostacoli} & torretta, proiettile, barra rotante \\
\texttt{Content/Blueprints/Livelli/Gym} & cresce con la distanza, cambia modello, presenza casuale \\
\texttt{Content/Blueprints/Livelli/Ambiente} & mulino, staccionata, spargi, luce tremolante, casse fisiche \\
\bottomrule
\end{tabular}
\end{center}

% =====================================================================
\section{Creare un livello da zero}

Questa è la procedura, uguale per i tre livelli. La si fa una volta per livello, poi si riempiono le tabelle una riga alla volta.

\subsection{Il file del livello}
\begin{enumerate}
  \item \nodo{File} $\rightarrow$ \nodo{New Level\dots} $\rightarrow$ \nodo{Empty Level} (completamente vuoto, senza luci: le mettiamo noi; non \emph{Empty Open World}, che crea un mondo a celle).
  \item \nodo{File} $\rightarrow$ \nodo{Save Current Level As\dots} $\rightarrow$ cartella \texttt{Maps}, nome \var{L\_GD13\_Livello1\_BP}.
  \item \nodo{Window} $\rightarrow$ \nodo{World Settings} $\rightarrow$ \emph{GameMode Override} = \var{GM\_Platform}. (Nei livelli C++ qui c'è \var{PlatformerGameMode}.)
\end{enumerate}

\subsection{Luce e atmosfera}

Ogni livello comincia con la sezione ``Luce, atmosfera e partenza'' della sua tabella. Gli actor si trovano in \nodo{Place Actors} (\tasto{Shift}+\tasto{1}) $\rightarrow$ \nodo{Lights} e \nodo{Visual Effects}, oppure tutti insieme con \nodo{Window} $\rightarrow$ \nodo{Env.~Light Mixer} $\rightarrow$ \emph{Create Sky Light}, \emph{Atmospheric Light}, \emph{Sky Atmosphere}, \emph{Height Fog}. Il \emph{Post Process Volume}: una volta piazzato, in Details cercate ``Unbound'' e spuntate \emph{Infinite Extent}; poi ogni valore (per esempio \emph{Bloom} $\rightarrow$ \emph{Intensity}) si attiva con la sua casella prima di cambiarlo. Infine \nodo{Player Start} (in \nodo{Place Actors} $\rightarrow$ \nodo{Basic}) alla posizione indicata.

\subsection{Le righe delle tabelle}

\begin{description}
  \item[Cubo, Pedana] \nodo{Cube}; Location = posizione, Scale = misure/100, materiale dalla colonna \emph{Impostazioni}. Se c'è \emph{Tag: WallRun}, aggiungete il tag.
  \item[Static Mesh] trascinate la mesh indicata dal Content Browser; Location, Rotation (yaw) e Scale come in tabella. ``Collision: No Collision'' si imposta in Details $\rightarrow$ \emph{Collision Presets}.
  \item[Text Render] \nodo{Place Actors} $\rightarrow$ cercate \emph{Text Render}. Nei Details: \emph{Text} (per andare a capo dentro il testo: \tasto{Shift}+\tasto{Invio}; nella tabella l'a capo è ``/''), \emph{World Size}, \emph{Text Render Color}, \emph{Horizontal Alignment} = Center, Rotation yaw 180 (così il testo guarda chi arriva da $-$X).
  \item[Decal] \nodo{Place Actors} $\rightarrow$ \emph{Decal Actor}; Rotation e \emph{Decal Size} come in tabella, \emph{Decal Material} indicato.
  \item[Point Light, Spot Light] Mobility \emph{Movable}; \emph{Intensity Units} = Candelas; Intensity, \emph{Light Color}, \emph{Attenuation Radius}; per lo spot anche \emph{Inner/Outer Cone Angle}; ``niente ombre'' = \emph{Cast Shadows} spento.
  \item[BP\_\dots] trascinate il vostro Blueprint nel livello, scrivete posizione e yaw, poi nei Details, sezione \emph{Default}, i valori delle variabili elencati in \emph{Impostazioni}. Perché compaiano nei Details le variabili devono essere \emph{Instance Editable} (l'occhio aperto accanto alla variabile nel Blueprint). Le variabili non nominate restano al loro default (catalogo, sezione~\ref{sec:catalogo}).
  \item[Riferimenti] \var{Targets = [Porta\_Doppia]}, \var{ListenTo = [\dots]}: è un array di actor. Nei Details, \nodo{+} per aggiungere un elemento, poi il contagocce e un clic sull'actor nel viewport (o sceglietelo dalla tendina).
\end{description}

\attenzione{Le zone (\var{BP\_Zona} e figli) si impostano con \var{ZoneExtent}, che è \emph{metà} della misura (come il \emph{Box Extent} di Unreal): \var{ZoneExtent} $(180;\ 180;\ 100)$ è una scatola di $3{,}6 \times 3{,}6 \times 2$\,m.}

\subsection{La scorciatoia: lo scheletro dei livelli con uno script}

Mettere a mano duecento righe è un ottimo esercizio di level design, ma è lungo. Se preferite concentrarvi sui Blueprint, lo script può costruire per voi la \emph{struttura} dei tre livelli \var{\_BP}: blocchi, luci, atmosfera, scritte, decal, piante e Player Start, \emph{senza} nessuna classe C++. Dove andrebbe un actor C++ mette invece un \textbf{segnaposto}: una scritta arancione, visibile solo nell'editor, con il nome del Blueprint da piazzare lì e il nome che deve avere nell'Outliner (cartella \texttt{DaFare}).

\begin{enumerate}
  \item \emph{Output Log} in modalità Python: \texttt{import gd13\_platformer as g; g.costruisci\_scheletro\_bp()}.
  \item In ogni livello \var{\_BP}: \emph{World Settings} $\rightarrow$ \emph{GameMode Override} = \var{GM\_Platform}.
  \item Per ogni segnaposto: piazzate il Blueprint indicato \emph{alla stessa posizione} (selezionate il segnaposto, copiate la \emph{Location} con tasto destro $\rightarrow$ \emph{Copy}, incollatela nel Blueprint), impostate le variabili dalla tabella, poi cancellate il segnaposto. Quando la cartella \texttt{DaFare} è vuota, il livello è finito.
\end{enumerate}

Lo script \textbf{non sovrascrive} un livello \var{\_BP} che esiste già: per rifarlo, cancellatelo prima a mano.

% =====================================================================
\section{Il catalogo dei Blueprint}\label{sec:catalogo}
\addtocontents{toc}{\protect\setcounter{tocdepth}{1}}

Questi sono tutti i Blueprint che compaiono nei tre livelli. Ognuno rifà una classe C++ e ha le \emph{stesse variabili con lo stesso nome} (senza la \var{b} iniziale dei booleani del C++): così le tabelle dei livelli valgono per tutte e due le versioni, e se un comportamento non vi torna potete aprire il file \texttt{.h} indicato e leggere il commento. La logica a nodi è nella guida Ninja, al capitolo indicato.

\subsection*{Quali servono, e dove}

%% generato da strumenti/estrai_livelli.py: non modificare a mano
\begin{longtable}{lcccl}
\toprule
Blueprint & Livello 1 & Livello 2 & Livello 3 & Sezione \\
\midrule\endhead
\texttt{BP\_Ascensore} & -- & 2 & -- & \ref{bp:BP_Ascensore} \\
\texttt{BP\_BarraRotante} & -- & 1 & 1 & \ref{bp:BP_BarraRotante} \\
\texttt{BP\_CambiaModello} & 1 & -- & 1 & \ref{bp:BP_CambiaModello} \\
\texttt{BP\_CasseFisiche} & 1 & -- & -- & \ref{bp:BP_CasseFisiche} \\
\texttt{BP\_CresceConDistanza} & 1 & 3 & -- & \ref{bp:BP_CresceConDistanza} \\
\texttt{BP\_Leva} & -- & 1 & -- & \ref{bp:BP_Leva} \\
\texttt{BP\_LuceAttivabile} & 1 & -- & -- & \ref{bp:BP_LuceAttivabile} \\
\texttt{BP\_LuceTremolante} & -- & -- & 9 & \ref{bp:BP_LuceTremolante} \\
\texttt{BP\_Moneta} & 10 & 10 & 5 & \ref{bp:BP_Moneta} \\
\texttt{BP\_Mulino} & 1 & -- & -- & \ref{bp:BP_Mulino} \\
\texttt{BP\_Negozio} & -- & 1 & -- & \ref{bp:BP_Negozio} \\
\texttt{BP\_PiastraPressione} & 1 & -- & -- & \ref{bp:BP_PiastraPressione} \\
\texttt{BP\_PiattaformaCrolla} & -- & 3 & -- & \ref{bp:BP_PiattaformaCrolla} \\
\texttt{BP\_PiattaformaMobile} & 1 & -- & 2 & \ref{bp:BP_PiattaformaMobile} \\
\texttt{BP\_PickupAbilita} & 2 & 1 & -- & \ref{bp:BP_PickupAbilita} \\
\texttt{BP\_PickupBatteria} & -- & -- & 2 & \ref{bp:BP_PickupBatteria} \\
\texttt{BP\_PickupSalute} & -- & 1 & 1 & \ref{bp:BP_PickupSalute} \\
\texttt{BP\_Porta} & 1 & -- & -- & \ref{bp:BP_Porta} \\
\texttt{BP\_PresenzaCasuale} & 1 & -- & -- & \ref{bp:BP_PresenzaCasuale} \\
\texttt{BP\_Spargi} & 2 & -- & -- & \ref{bp:BP_Spargi} \\
\texttt{BP\_Staccionata} & 1 & -- & -- & \ref{bp:BP_Staccionata} \\
\texttt{BP\_Teletrasporto} & -- & -- & 1 & \ref{bp:BP_Teletrasporto} \\
\texttt{BP\_TimerLivello} & -- & -- & 1 & \ref{bp:BP_TimerLivello} \\
\texttt{BP\_Torretta} & -- & 2 & 1 & \ref{bp:BP_Torretta} \\
\texttt{BP\_Traguardo} & 1 & 1 & 1 & \ref{bp:BP_Traguardo} \\
\texttt{BP\_Trampolino} & -- & 2 & -- & \ref{bp:BP_Trampolino} \\
\texttt{BP\_ZonaCaduta} & 1 & 1 & 1 & \ref{bp:BP_ZonaCaduta} \\
\texttt{BP\_ZonaDanno} & -- & 1 & -- & \ref{bp:BP_ZonaDanno} \\
\texttt{BP\_ZonaModificatore} & 1 & 1 & 1 & \ref{bp:BP_ZonaModificatore} \\
\bottomrule
\end{longtable}


\textbf{Ordine consigliato.} Prima i quattro padri (\var{BP\_Zona}, \var{BP\_Pickup}, \var{BP\_Interruttore}, \var{BP\_Attivabile}), poi \var{BP\_ZonaCaduta}, \var{BP\_Moneta}, \var{BP\_Traguardo} e \var{BP\_PiattaformaMobile}: con questi il livello 1 si gioca dall'inizio alla fine. Il resto si aggiunge tratto per tratto. I padri si creano con \nodo{Blueprint Class} $\rightarrow$ \nodo{Actor}; i figli con tasto destro sul padre $\rightarrow$ \nodo{Create Child Blueprint Class}.

%% generato da strumenti/estrai_livelli.py: non modificare a mano
\subsection{\texttt{BP\_Zona}}\label{bp:BP_Zona}
\noindent Rifà \texttt{APlatformerZone} (\texttt{PlatformerZone.h}). Classe padre: \nodo{Actor}.\par
\smallskip\noindent\textbf{Componenti.} \nodo{Box Collision} \var{Trigger} (root; Collision Preset \emph{Trigger}), \nodo{Static Mesh} \var{Visual} (un \var{Cube}, senza collisione), \nodo{Text Render} \var{Scritta}.\par
\noindent\textbf{Origine (pivot).} il \textbf{centro} della scatola. Nel \nodo{Construction Script}: \nodo{Set Box Extent} = \var{ZoneExtent}; \var{Visual} con scala $(\mathrm{ZoneExtent.X}/50;\ \mathrm{ZoneExtent.Y}/50;\ \mathrm{VisualThickness}/100)$ e Z relativa $-\mathrm{ZoneExtent.Z}+\mathrm{VisualThickness}/2$ (sta sul fondo); \nodo{Set Visibility} di \var{Visual} = \var{ShowVisual}; la scritta a Z $\mathrm{ZoneExtent.Z}+20$.\par
\smallskip\noindent\textbf{Variabili} (tutte \emph{Instance Editable}):
\begin{variabili}
ZoneExtent & Vector & (150, 150, 100) & Metà delle dimensioni della zona, in cm (come la Box Extent di un Box Collision). \\
OnlyPlayer & Boolean & sì & Se vero, reagisce solo al personaggio del giocatore (non ad altri actor). \\
ShowVisual & Boolean & sì & Mostra una pedana sul fondo della zona, così la si vede anche in gioco. \\
VisualThickness & Float & 4 & Spessore della pedana visibile sul fondo della zona, in cm. \\
VisualMaterial & MaterialInterface (riferimento) &  & Materiale della pedana visibile (opzionale). \\
Label & Text &  & Scritta sopra la zona (vuota = nessuna scritta). \\
\end{variabili}
\smallskip\noindent\textbf{Logica.} filtra il giocatore e chiama gli eventi vuoti \var{Entra}/\var{Esce} che i figli riempiono (Ninja, cap.~10).\par

\subsection{\texttt{BP\_Pickup}}\label{bp:BP_Pickup}
\noindent Rifà \texttt{APlatformerPickup} (\texttt{PlatformerPickups.h}). Classe padre: \nodo{Actor}.\par
\smallskip\noindent\textbf{Componenti.} \nodo{Sphere Collision} \var{Sfera} (root, raggio \var{PickupRadius}), \nodo{Static Mesh} \var{Mesh} (senza collisione), \nodo{Point Light} \var{Luce} (intensità \var{GlowIntensity}, colore \var{GlowColor}, raggio 200), \nodo{Text Render} \var{Scritta} a Z 90.\par
\noindent\textbf{Origine (pivot).} il \textbf{centro} dell'oggetto (che gira e ondeggia intorno a quel punto).\par
\smallskip\noindent\textbf{Variabili} (tutte \emph{Instance Editable}):
\begin{variabili}
SpinSpeed & Float & 120 & Gradi al secondo di rotazione. \\
BobHeight & Float & 12 & Ampiezza dell'ondeggiamento verticale (cm). \\
BobSpeed & Float & 0.8 & Velocità dell'ondeggiamento (cicli al secondo). \\
PickupRadius & Float & 60 & Raggio di raccolta (cm). \\
Material & MaterialInterface (riferimento) &  &  \\
GlowIntensity & Float & 8 & Luce piccola che fa "brillare" l'oggetto e lo fa notare (0 = spenta). \\
GlowColor & Linear Color & (1, 0.8, 0.3) &  \\
PickupSound & SoundBase (riferimento) &  &  \\
Label & Text &  & Scritta sopra l'oggetto (vuota = nessuna). \\
\end{variabili}
\smallskip\noindent\textbf{Logica.} gira e ondeggia nel Tick; all'overlap chiama \var{Raccolto} e, se torna vero, suono e \nodo{Destroy Actor} (Ninja, cap.~11).\par

\subsection{\texttt{BP\_Interruttore}}\label{bp:BP_Interruttore}
\noindent Rifà \texttt{ASwitchBase} (\texttt{Switches.h}). Classe padre: \nodo{Actor}.\par
\smallskip\noindent\textbf{Componenti.} dipendono dal figlio (piastra o leva).\par
\smallskip\noindent\textbf{Variabili} (tutte \emph{Instance Editable}):
\begin{variabili}
Targets & Array di Activable &  & Gli attivabili da chiamare (primo metodo: For Each Loop sull'array). \\
ClickSound & SoundBase (riferimento) &  & Suono quando cambia stato (opzionale). \\
OnMaterial & MaterialInterface (riferimento) &  & Materiale "acceso" e "spento" della parte che si muove (opzionali). \\
OffMaterial & MaterialInterface (riferimento) &  &  \\
\end{variabili}
\smallskip\noindent\textbf{Logica.} \var{ImpostaStato}: chiama \var{Attiva}/\var{Disattiva} dell'interfaccia \var{BPI\_Attivabile} su ogni elemento di \var{Targets} e lancia il dispatcher \var{OnCambiato} (Ninja, cap.~9).\par

\subsection{\texttt{BP\_Attivabile}}\label{bp:BP_Attivabile}
\noindent Rifà \texttt{AActivable} (\texttt{Activable.h}). Classe padre: \nodo{Actor}.\par
\smallskip\noindent\textbf{Componenti.} una \nodo{Scene} \var{Root}; il resto lo aggiunge il figlio. Implementa l'interfaccia \var{BPI\_Attivabile}.\par
\smallskip\noindent\textbf{Variabili} (tutte \emph{Instance Editable}):
\begin{variabili}
StartActive & Boolean & no & Stato all'avvio. \\
ListenTo & Array di SwitchBase &  & Interruttori di cui ascoltare l'event dispatcher (secondo metodo della lezione). \\
TransitionTime & Float & 1 & Secondi della "timeline" di attivazione (porta che si apre, piattaforma che sale...). \\
TransitionCurve & CurveFloat (riferimento) &  & Curva della timeline (opzionale): senza curva il movimento è morbido (ease in/out). \\
ActivateSound & SoundBase (riferimento) &  & Suono all'attivazione (opzionale). \\
\end{variabili}
\smallskip\noindent\textbf{Logica.} \var{Attiva} $\rightarrow$ timeline lunga \var{TransitionTime} in avanti, \var{Disattiva} $\rightarrow$ indietro; l'\emph{Update} chiama la funzione vuota \var{Applica(Alpha)} che i figli ridefiniscono. In BeginPlay: per ogni interruttore in \var{ListenTo}, \nodo{Bind Event to OnCambiato}. Se \var{StartActive}, parte già attivo (Ninja, cap.~9).\par

\subsection{\texttt{BP\_ZonaCaduta}}\label{bp:BP_ZonaCaduta}
\noindent Rifà \texttt{AFallZone} (\texttt{PlatformerZones.h}). Classe padre: \nodo{BP\_Zona}.\par
\noindent In più eredita le variabili di \texttt{BP\_Zona} (sezione~\ref{bp:BP_Zona}).\par
\smallskip\noindent\textbf{Logica.} \var{Entra} $\rightarrow$ \var{GestisciCaduta} del componente \var{AC\_PostoSicuro}: perde una vita e torna all'ultimo posto sicuro; a zero vite, game over (Ninja, cap.~5). Nei livelli è invisibile (\var{ShowVisual} = no).\par

\subsection{\texttt{BP\_ZonaMortale}}\label{bp:BP_ZonaMortale}
\noindent Rifà \texttt{AKillZone} (\texttt{PlatformerZones.h}). Classe padre: \nodo{BP\_Zona}.\par
\smallskip\noindent\textbf{Variabili} (tutte \emph{Instance Editable}):
\begin{variabili}
RestartDelay & Float & 0.3 & Secondi prima che il livello riparta. \\
\end{variabili}
\noindent In più eredita le variabili di \texttt{BP\_Zona} (sezione~\ref{bp:BP_Zona}).\par
\smallskip\noindent\textbf{Logica.} \var{Entra} $\rightarrow$ \var{UccidiERiavvia}(\var{RestartDelay}) (Ninja, cap.~8).\par

\subsection{\texttt{BP\_ZonaDanno}}\label{bp:BP_ZonaDanno}
\noindent Rifà \texttt{ADamageZone} (\texttt{PlatformerZones.h}). Classe padre: \nodo{BP\_Zona}.\par
\smallskip\noindent\textbf{Variabili} (tutte \emph{Instance Editable}):
\begin{variabili}
DamagePerHit & Float & 20 & Danno inflitto a ogni "colpo". \\
HitInterval & Float & 0.5 & Secondi tra un colpo e l'altro, finché si resta nella zona. \\
\end{variabili}
\noindent In più eredita le variabili di \texttt{BP\_Zona} (sezione~\ref{bp:BP_Zona}).\par
\smallskip\noindent\textbf{Logica.} timer ripetuto ogni \var{HitInterval} che fa \nodo{Apply Damage} di \var{DamagePerHit} finché si è dentro (Ninja, cap.~8, la lava).\par

\subsection{\texttt{BP\_ZonaModificatore}}\label{bp:BP_ZonaModificatore}
\noindent Rifà \texttt{AModifierZone} (\texttt{PlatformerZones.h}). Classe padre: \nodo{BP\_Zona}.\par
\smallskip\noindent\textbf{Variabili} (tutte \emph{Instance Editable}):
\begin{variabili}
WalkSpeedMultiplier & Float & 1 & Moltiplica la velocità di camminata (1 = invariata). \\
JumpMultiplier & Float & 1 & Moltiplica la forza del salto (1 = invariata). \\
GravityMultiplier & Float & 1 & Moltiplica la gravità (meno di 1 = salti "lunari"). \\
GroundFrictionMultiplier & Float & 1 & Moltiplica l'attrito a terra (Ground Friction). Meno di 1 = si scivola: il ghiaccio. \\
BrakingMultiplier & Float & 1 & Moltiplica la frenata (Braking Deceleration Walking). Meno di 1 = ci si ferma tardi. \\
DurationAfterExit & Float & 3 & Quanto dura l'effetto dopo essere USCITI dalla zona, in secondi. 0 = finisce appena si esce. Negativo = dura per tutto il livello. \\
\end{variabili}
\noindent In più eredita le variabili di \texttt{BP\_Zona} (sezione~\ref{bp:BP_Zona}).\par
\smallskip\noindent\textbf{Logica.} all'ingresso salva e moltiplica velocità, salto, gravità, attrito e frenata del Character Movement; all'uscita li ripristina dopo \var{DurationAfterExit} secondi (Ninja, cap.~10).\par

\subsection{\texttt{BP\_Teletrasporto}}\label{bp:BP_Teletrasporto}
\noindent Rifà \texttt{ATeleportZone} (\texttt{PlatformerZones.h}). Classe padre: \nodo{BP\_Zona}.\par
\smallskip\noindent\textbf{Variabili} (tutte \emph{Instance Editable}):
\begin{variabili}
Destination & Vector & (500, 0, 100) & Punto di arrivo, relativo alla zona (si trascina nel viewport con il rombo). \\
StopVelocity & Boolean & sì & Se vero, all'arrivo si azzera la velocità. \\
\end{variabili}
\noindent In più eredita le variabili di \texttt{BP\_Zona} (sezione~\ref{bp:BP_Zona}).\par
\smallskip\noindent\textbf{Logica.} \nodo{Teleport} dell'eroe in \nodo{Transform Location}(\var{Destination}) (la destinazione è \emph{relativa} alla zona: spuntate \emph{Show 3D Widget} sulla variabile per trascinarla nel viewport); se \var{StopVelocity}, velocità a zero (Ninja, cap.~10).\par

\subsection{\texttt{BP\_Negozio}}\label{bp:BP_Negozio}
\noindent Rifà \texttt{AShopZone} (\texttt{PlatformerZones.h}). Classe padre: \nodo{BP\_Zona}.\par
\smallskip\noindent\textbf{Variabili} (tutte \emph{Instance Editable}):
\begin{variabili}
Cost & Integer & 5 & Monete richieste. \\
Targets & Array di Activable &  & Cosa si attiva pagando (porte, piattaforme, luci...). \\
OnlyOnce & Boolean & sì & Se vero, si paga una volta sola. \\
\end{variabili}
\noindent In più eredita le variabili di \texttt{BP\_Zona} (sezione~\ref{bp:BP_Zona}).\par
\smallskip\noindent\textbf{Logica.} se le monete bastano, le toglie (\var{Cost}) e chiama \var{Attiva} su ogni elemento di \var{Targets}; con \var{OnlyOnce} funziona una volta sola. Altrimenti un messaggio ``Servono N monete'' (Ninja, cap.~10).\par

\subsection{\texttt{BP\_Traguardo}}\label{bp:BP_Traguardo}
\noindent Rifà \texttt{ALevelGoal} (\texttt{PlatformerZones.h}). Classe padre: \nodo{BP\_Zona}.\par
\smallskip\noindent\textbf{Variabili} (tutte \emph{Instance Editable}):
\begin{variabili}
NextLevel & Name &  & Nome della mappa da aprire, per esempio L\_GD13\_Livello2. \\
Message & String & Livello completato! & Messaggio a schermo quando si arriva. \\
Delay & Float & 1 & Secondi prima del cambio di livello. \\
ResetDeathCount & Boolean & no & Se vero, azzera il contatore delle morti (da usare sull'ultimo traguardo). \\
\end{variabili}
\noindent In più eredita le variabili di \texttt{BP\_Zona} (sezione~\ref{bp:BP_Zona}).\par
\smallskip\noindent\textbf{Logica.} \nodo{Do Once} $\rightarrow$ incassa le monete, registra il tempo, mostra \var{Message}, se \var{ResetDeathCount} azzera le morti, \nodo{Delay} \var{Delay} $\rightarrow$ \var{VaiAlLivello(NextLevel)} (Ninja, cap.~10). Nei livelli \_BP \var{NextLevel} è il livello \_BP successivo; dall'ultimo si torna a \var{L\_GD13\_Menu}.\par

\subsection{\texttt{BP\_TimerLivello}}\label{bp:BP_TimerLivello}
\noindent Rifà \texttt{ALevelTimer} (\texttt{LevelTimer.h}). Classe padre: \nodo{Actor}.\par
\smallskip\noindent\textbf{Componenti.} \nodo{Scene} root, \nodo{Text Render} \var{Testo} (World Size 120, centrato).\par
\noindent\textbf{Origine (pivot).} dove deve stare la scritta del tempo.\par
\smallskip\noindent\textbf{Variabili} (tutte \emph{Instance Editable}):
\begin{variabili}
TimeLimit & Float & 60 & Secondi a disposizione per finire il livello. \\
WarningThreshold & Float & 10 & Sotto questa soglia il tempo diventa rosso. \\
ShowOnScreen & Boolean & no & Scrive il tempo anche come messaggio di debug in alto a sinistra (l'HUD lo mostra già). \\
ShowInWorld & Boolean & sì & Scrive il tempo anche nel mondo, sulla scritta dell'actor (utile se lo si piazza in vista). \\
\end{variabili}
\smallskip\noindent\textbf{Logica.} timer ogni 0,1\,s che scala \var{TimeLimit}; sotto \var{WarningThreshold} il testo diventa rosso; a zero, \var{UccidiERiavvia} (Ninja, cap.~10). \var{ShowOnScreen} = mostra il tempo anche nell'HUD.\par

\subsection{\texttt{BP\_Moneta}}\label{bp:BP_Moneta}
\noindent Rifà \texttt{ACoinPickup} (\texttt{PlatformerPickups.h}). Classe padre: \nodo{BP\_Pickup}.\par
\smallskip\noindent\textbf{Componenti.} quelli di \var{BP\_Pickup}; la mesh è un \var{Cylinder} schiacciato (scala $0{,}45 \times 0{,}45 \times 0{,}08$) ruotato di 90\textdegree{} in Roll, così la si vede di taglio.\par
\smallskip\noindent\textbf{Variabili} (tutte \emph{Instance Editable}):
\begin{variabili}
Value & Integer & 1 &  \\
\end{variabili}
\noindent In più eredita le variabili di \texttt{BP\_Pickup} (sezione~\ref{bp:BP_Pickup}).\par
\smallskip\noindent\textbf{Logica.} \var{Raccolto}: \var{GI.Monete} += \var{Value}, ritorna vero (Ninja, cap.~11).\par

\subsection{\texttt{BP\_PickupAbilita}}\label{bp:BP_PickupAbilita}
\noindent Rifà \texttt{AAbilityPickup} (\texttt{PlatformerPickups.h}). Classe padre: \nodo{BP\_Pickup}.\par
\smallskip\noindent\textbf{Componenti.} quelli di \var{BP\_Pickup}; mesh \var{Sphere} scala 0,4.\par
\smallskip\noindent\textbf{Variabili} (tutte \emph{Instance Editable}):
\begin{variabili}
Ability & E\_Abilita & DoubleJump &  \\
Message & String & Doppio salto sbloccato! Premi salto di nuovo in aria & Messaggio a schermo quando la si prende. \\
HideIfAlreadyUnlocked & Boolean & sì & Se vero, sparisce quando l'abilità è già stata sbloccata (per esempio rigiocando il livello). \\
\end{variabili}
\noindent In più eredita le variabili di \texttt{BP\_Pickup} (sezione~\ref{bp:BP_Pickup}).\par
\smallskip\noindent\textbf{Logica.} \var{Raccolto}: \var{SbloccaAbilita(Ability)}, \var{MostraMessaggio(Message)}. Se \var{HideIfAlreadyUnlocked} e l'abilità c'è già, in BeginPlay si distrugge. Il cartello si scrive da solo nel Construction Script (Ninja, cap.~11).\par

\subsection{\texttt{BP\_PickupSalute}}\label{bp:BP_PickupSalute}
\noindent Rifà \texttt{AHealthPickup} (\texttt{PlatformerPickups.h}). Classe padre: \nodo{BP\_Pickup}.\par
\smallskip\noindent\textbf{Componenti.} quelli di \var{BP\_Pickup}; mesh \var{Cylinder} scala $(0{,}6;\ 0{,}6;\ 0{,}9)$, materiale verde.\par
\smallskip\noindent\textbf{Variabili} (tutte \emph{Instance Editable}):
\begin{variabili}
HealAmount & Float & 40 &  \\
\end{variabili}
\noindent In più eredita le variabili di \texttt{BP\_Pickup} (sezione~\ref{bp:BP_Pickup}).\par
\smallskip\noindent\textbf{Logica.} \var{Raccolto}: se la vita non è piena, \nodo{Apply Damage} di $-$\var{HealAmount} all'eroe e ritorna vero; altrimenti falso e resta lì (Ninja, cap.~11).\par

\subsection{\texttt{BP\_PickupBatteria}}\label{bp:BP_PickupBatteria}
\noindent Rifà \texttt{ABatteryPickup} (\texttt{PlatformerPickups.h}). Classe padre: \nodo{BP\_Pickup}.\par
\smallskip\noindent\textbf{Componenti.} quelli di \var{BP\_Pickup}; mesh \var{Cylinder} scala $(0{,}25;\ 0{,}25;\ 0{,}45)$.\par
\smallskip\noindent\textbf{Variabili} (tutte \emph{Instance Editable}):
\begin{variabili}
BatterySeconds & Float & 15 & Secondi di luce aggiunti. \\
\end{variabili}
\noindent In più eredita le variabili di \texttt{BP\_Pickup} (sezione~\ref{bp:BP_Pickup}).\par
\smallskip\noindent\textbf{Logica.} \var{Raccolto}: \var{Batteria} della torcia += \var{BatterySeconds} (con \nodo{Min} sul massimo) (Ninja, cap.~4 e 11).\par

\subsection{\texttt{BP\_PiattaformaMobile}}\label{bp:BP_PiattaformaMobile}
\noindent Rifà \texttt{AMovingPlatform} (\texttt{MovingPlatform.h}). Classe padre: \nodo{Actor}.\par
\smallskip\noindent\textbf{Componenti.} \nodo{Static Mesh} \var{Mesh} (root, \var{Cube} con scala \var{PlatformSize}/100, collisione \emph{BlockAll}), \nodo{Audio} \var{Suono} (\var{MoveSound}).\par
\noindent\textbf{Origine (pivot).} il \textbf{centro} della piattaforma (nelle tabelle la Z è già il centro: piano $-$ metà spessore).\par
\smallskip\noindent\textbf{Variabili} (tutte \emph{Instance Editable}):
\begin{variabili}
EndOffset & Vector & (0, 400, 0) & Spostamento del punto di arrivo rispetto alla posizione iniziale (si trascina col widget nel viewport). \\
TravelTime & Float & 2 & Secondi per andare da un estremo all'altro. \\
PauseTime & Float & 0.5 & Secondi di sosta a ogni estremo. \\
StartDelay & Float & 0 & Ritardo iniziale: sfasa piattaforme uguali messe in fila. \\
Smooth & Boolean & sì & Movimento morbido agli estremi (ease in/out) invece che lineare. \\
PlatformSize & Vector & (300, 300, 30) & Dimensioni della piattaforma in cm. \\
PlatformMaterial & MaterialInterface (riferimento) &  & Materiale della piattaforma (opzionale). \\
MoveSound & SoundBase (riferimento) &  & UE6: il suono della piattaforma mentre si muove (Sound Cue, opzionale). \\
PitchForward & Float & 1.15 & UE6: il pitch cambia con la direzione: all'andata... \\
PitchBackward & Float & 0.85 & ...e al ritorno. \\
\end{variabili}
\smallskip\noindent\textbf{Logica.} timeline di \var{TravelTime} secondi fra la posizione iniziale e iniziale + \var{EndOffset}, pausa \var{PauseTime} agli estremi, ritardo iniziale \var{StartDelay}; \nodo{Set Pitch Multiplier} del suono a \var{PitchForward} all'andata e \var{PitchBackward} al ritorno (Ninja, cap.~8).\par

\subsection{\texttt{BP\_PiattaformaCrolla}}\label{bp:BP_PiattaformaCrolla}
\noindent Rifà \texttt{ACrumblingPlatform} (\texttt{PlatformerProps.h}). Classe padre: \nodo{Actor}.\par
\smallskip\noindent\textbf{Componenti.} \nodo{Static Mesh} \var{Mesh} (root, \var{Cube} scala \var{PlatformSize}/100), \nodo{Box Collision} \var{Sopra} (extent $(\mathrm{X}/2-5;\ \mathrm{Y}/2-5;\ 20)$, a Z $\mathrm{PlatformSize.Z}/2+20$: sente chi ci sale).\par
\noindent\textbf{Origine (pivot).} il \textbf{centro} della piattaforma.\par
\smallskip\noindent\textbf{Variabili} (tutte \emph{Instance Editable}):
\begin{variabili}
PlatformSize & Vector & (250, 250, 30) &  \\
ShakeTime & Float & 0.8 & Secondi in cui trema prima di cadere. \\
ShakeAmount & Float & 4 & Quanto trema (cm). \\
RespawnTime & Float & 3 & Secondi prima di ricomparire (0 = non ricompare). \\
PlatformMaterial & MaterialInterface (riferimento) &  &  \\
WarningMaterial & MaterialInterface (riferimento) &  & Materiale mentre trema (avviso visivo, opzionale). \\
CrumbleSound & SoundBase (riferimento) &  &  \\
\end{variabili}
\smallskip\noindent\textbf{Logica.} qualcuno sale $\rightarrow$ materiale \var{WarningMaterial} e tremolio per \var{ShakeTime} $\rightarrow$ suono, nascosta e senza collisione $\rightarrow$ dopo \var{RespawnTime} torna (Ninja, cap.~8).\par

\subsection{\texttt{BP\_Trampolino}}\label{bp:BP_Trampolino}
\noindent Rifà \texttt{AJumpPad} (\texttt{PlatformerProps.h}). Classe padre: \nodo{Actor}.\par
\smallskip\noindent\textbf{Componenti.} \var{Base} (\var{Cylinder} scala $(1{,}4;\ 1{,}4;\ 0{,}2)$), \var{Pad} (\var{Cylinder} scala $(0{,}85;\ 0{,}85;\ 0{,}4)$ a Z 60, materiale \var{PadMaterial}), \nodo{Box Collision} \var{Trigger} extent $(40;\ 40;\ 150)$ a Z 250. \emph{Attenzione}: il trigger sta sopra il pad, così il lancio scatta mentre si passa sopra.\par
\noindent\textbf{Origine (pivot).} \textbf{a terra}, al centro della base.\par
\smallskip\noindent\textbf{Variabili} (tutte \emph{Instance Editable}):
\begin{variabili}
LaunchVelocity & Vector & (300, 0, 1300) & Velocità data al personaggio, nello spazio dell'actor (X avanti, Z su). cm/s. \\
OverrideHorizontal & Boolean & sì & Se vero, sostituisce la velocità orizzontale invece di sommarla. \\
BounceSound & SoundBase (riferimento) &  &  \\
PadMaterial & MaterialInterface (riferimento) &  &  \\
\end{variabili}
\smallskip\noindent\textbf{Logica.} overlap $\rightarrow$ \nodo{Launch Character} con \nodo{Rotate Vector}(\var{LaunchVelocity}, rotazione dell'actor), \emph{XY Override} = \var{OverrideHorizontal}, \emph{Z Override} sì; il pad si schiaccia per un attimo (Ninja, cap.~8).\par

\subsection{\texttt{BP\_BarraRotante}}\label{bp:BP_BarraRotante}
\noindent Rifà \texttt{ASpinningObstacle} (\texttt{SpinningObstacle.h}). Classe padre: \nodo{Actor}.\par
\smallskip\noindent\textbf{Componenti.} \var{Base} (\var{Cylinder} scala $(0{,}4;\ 0{,}4;\ \mathrm{BarHeight}/100)$ a Z $\mathrm{BarHeight}/2$), \var{Perno} (Scene), sotto il perno \var{Barra} (\var{Cube} scala \var{BarSize}/100 a Z \var{BarHeight}, collisione \emph{OverlapAllDynamic}: \textbf{non} blocca).\par
\noindent\textbf{Origine (pivot).} \textbf{a terra}, al centro.\par
\smallskip\noindent\textbf{Variabili} (tutte \emph{Instance Editable}):
\begin{variabili}
DegreesPerSecond & Float & 90 & Gradi al secondo attorno all'asse verticale (negativo = senso opposto). \\
BarSize & Vector & (600, 30, 30) & Lunghezza, spessore e altezza della sbarra, in cm. \\
BarHeight & Float & 60 & Altezza della sbarra dal pavimento, in cm. \\
PushStrength & Float & 900 & Forza con cui la sbarra spinge via il personaggio. \\
KillOnTouch & Boolean & no & Se vero, toccare la sbarra fa ripartire il livello invece di spingere. \\
\end{variabili}
\smallskip\noindent\textbf{Logica.} \nodo{Rotating Movement} sul perno a \var{DegreesPerSecond}; all'overlap della barra: se \var{KillOnTouch} $\rightarrow$ \var{UccidiERiavvia}(0,3), altrimenti \nodo{Launch Character} lontano dal centro con forza \var{PushStrength} (Ninja, cap.~8).\par

\subsection{\texttt{BP\_Torretta}}\label{bp:BP_Torretta}
\noindent Rifà \texttt{ATurret} (\texttt{Turret.h}). Classe padre: \nodo{Actor}.\par
\smallskip\noindent\textbf{Componenti.} \var{Base} (\var{Cylinder} scala $(0{,}8;\ 0{,}8;\ 1)$ a Z 50), \var{Testa} (Scene a Z 130) con \var{TestaMesh} (\var{Sphere} 0,75), \var{Canna} (\var{Cylinder} $0{,}22 \times 0{,}22 \times 1$ ruotato in avanti) e \var{Bocca} (Scene a X 110); \nodo{Sphere Collision} \var{Rilevamento} raggio \var{Range}.\par
\noindent\textbf{Origine (pivot).} \textbf{a terra} sopra la colonna rossa; yaw = verso cui guarda a riposo.\par
\smallskip\noindent\textbf{Variabili} (tutte \emph{Instance Editable}):
\begin{variabili}
Range & Float & 1800 & Raggio d'azione (cm). \\
FireInterval & Float & 1.4 & Secondi fra un colpo e l'altro. \\
TurnSpeed & Float & 3 & Velocità di rotazione della testa (RInterp To). \\
MinPitch & Float & -40 & Clamp del pitch della canna (gradi). \\
MaxPitch & Float & 30 &  \\
FireConeDegrees & Float & 12 & Spara solo se la canna è entro questo angolo dal bersaglio (gradi). \\
AimHeightOffset & Float & -60 & Dove mirare rispetto al centro del personaggio: negativo = ai piedi (cm). \\
LeadTarget & Boolean & no & Se vero, mira dove il personaggio SARÀ (anticipa il movimento). \\
ProjectileSpeed & Float & 1500 &  \\
ProjectileDamage & Float & 25 &  \\
ProjectileClass & Classe di PlatformerProjectile &  &  \\
RestartLevelOnHit & Boolean & no & Se vero, un proiettile che colpisce il personaggio fa ripartire il livello. \\
BodyMaterial & MaterialInterface (riferimento) &  &  \\
ProjectileMaterial & MaterialInterface (riferimento) &  &  \\
FireSound & SoundBase (riferimento) &  &  \\
\end{variabili}
\smallskip\noindent\textbf{Logica.} Ninja, cap.~8: \nodo{Find Look at Rotation} + \nodo{RInterp To} (\var{TurnSpeed}), pitch fra \var{MinPitch} e \var{MaxPitch}, mira all'altezza del bersaglio + \var{AimHeightOffset}; spara ogni \var{FireInterval} se l'angolo è entro \var{FireConeDegrees} e la vista è libera. Con \var{LeadTarget} mira dove il giocatore \emph{sarà}: posizione + velocità $\times$ (distanza / \var{ProjectileSpeed}). Ai proiettili passa \var{ProjectileSpeed}, \var{ProjectileDamage} e \var{RestartLevelOnHit}.\par

\subsection{\texttt{BP\_PiastraPressione}}\label{bp:BP_PiastraPressione}
\noindent Rifà \texttt{APressurePlate} (\texttt{Switches.h}). Classe padre: \nodo{BP\_Interruttore}.\par
\smallskip\noindent\textbf{Componenti.} \var{Cornice} (\var{Cube} scala $(1{,}6;\ 1{,}6;\ 0{,}08)$ a Z 4), \var{Piastra} (\var{Cube} $(1{,}3;\ 1{,}3;\ 0{,}1)$ a Z 10), \nodo{Box Collision} extent $(65;\ 65;\ 40)$ a Z 40.\par
\noindent\textbf{Origine (pivot).} \textbf{a terra}, al centro.\par
\smallskip\noindent\textbf{Variabili} (tutte \emph{Instance Editable}):
\begin{variabili}
StayPressed & Boolean & sì & Se vero, resta premuto anche dopo che si scende. \\
PressDepth & Float & 8 & Quanto si abbassa la pedana (cm). \\
PressTime & Float & 0.25 & Secondi della "timeline" di discesa. \\
\end{variabili}
\noindent In più eredita le variabili di \texttt{BP\_Interruttore} (sezione~\ref{bp:BP_Interruttore}).\par
\smallskip\noindent\textbf{Logica.} un Character entra $\rightarrow$ \var{ImpostaStato}(vero), piastra giù di \var{PressDepth} in \var{PressTime}, materiale \var{OnMaterial}; se \var{StayPressed} è no, quando esce torna su (Ninja, cap.~9).\par

\subsection{\texttt{BP\_Leva}}\label{bp:BP_Leva}
\noindent Rifà \texttt{AInteractLever} (\texttt{Switches.h}). Classe padre: \nodo{BP\_Interruttore}.\par
\smallskip\noindent\textbf{Componenti.} \var{Base} (\var{Cube} $(0{,}5;\ 0{,}7;\ 0{,}9)$ a Z 45), \var{Perno} (Scene a Z 90) con \var{Manico} (\var{Cylinder} $0{,}08 \times 0{,}08 \times 0{,}7$ a Z 35), \nodo{Text Render} a Z 190. Collisione della base sul canale \emph{Interact} (o \emph{Visibility}). Implementa \var{BPI\_Interagibile}.\par
\noindent\textbf{Origine (pivot).} \textbf{a terra}.\par
\smallskip\noindent\textbf{Variabili} (tutte \emph{Instance Editable}):
\begin{variabili}
Prompt & Text & Tira la leva & Testo mostrato dall'HUD vicino alla leva. \\
Toggle & Boolean & sì & Se vero, ogni pressione cambia stato; se falso si tira una volta sola. \\
Label & Text &  & Scritta sopra la leva (vuota = nessuna). \\
\end{variabili}
\noindent In più eredita le variabili di \texttt{BP\_Interruttore} (sezione~\ref{bp:BP_Interruttore}).\par
\smallskip\noindent\textbf{Logica.} \var{Interagisci}: se \var{Toggle} alterna acceso/spento, altrimenti accende soltanto; il manico ruota da $-$40\textdegree{} a +40\textdegree{}; \var{GetPrompt} restituisce \var{Prompt} (Ninja, cap.~9).\par

\subsection{\texttt{BP\_Porta}}\label{bp:BP_Porta}
\noindent Rifà \texttt{AActivableDoor} (\texttt{Activable.h}). Classe padre: \nodo{BP\_Attivabile}.\par
\smallskip\noindent\textbf{Componenti.} due \nodo{Scene} cerniere \var{CernieraSx} (Y $= -\mathrm{DoorWidth}/2$) e \var{CernieraDx} (Y $= +\mathrm{DoorWidth}/2$), con figlie le ante (\var{Cube} scala $(\mathrm{DoorThickness};\ \mathrm{DoorWidth}/2;\ \mathrm{DoorHeight})/100$, a Y $\pm\mathrm{DoorWidth}/4$ verso il centro e Z $\mathrm{DoorHeight}/2$).\par
\noindent\textbf{Origine (pivot).} \textbf{a terra}, al centro del vano.\par
\smallskip\noindent\textbf{Variabili} (tutte \emph{Instance Editable}):
\begin{variabili}
DoorWidth & Float & 240 & Larghezza totale del vano (le due ante insieme), in cm. \\
DoorHeight & Float & 300 &  \\
DoorThickness & Float & 20 &  \\
Sliding & Boolean & sì & Vero: le ante scorrono di lato. Falso: ruotano sui cardini. \\
OpenAngle & Float & 95 & Gradi di apertura se ruotano. \\
DoorMaterial & MaterialInterface (riferimento) &  &  \\
\end{variabili}
\noindent In più eredita le variabili di \texttt{BP\_Attivabile} (sezione~\ref{bp:BP_Attivabile}).\par
\smallskip\noindent\textbf{Logica.} \var{Applica(Alpha)}: se \var{Sliding}, le cerniere scorrono di $\mathrm{DoorWidth}/2 \cdot$ Alpha verso l'esterno; altrimenti ruotano di $\pm\mathrm{OpenAngle} \cdot$ Alpha (Ninja, cap.~9).\par

\subsection{\texttt{BP\_LuceAttivabile}}\label{bp:BP_LuceAttivabile}
\noindent Rifà \texttt{AActivableLight} (\texttt{Activable.h}). Classe padre: \nodo{BP\_Attivabile}.\par
\smallskip\noindent\textbf{Componenti.} \nodo{Point Light} (40 cd, Movable) e, se \var{ShowBulb}, una \var{Sphere} scala 0,2 con un materiale emissivo.\par
\noindent\textbf{Origine (pivot).} dove sta la lampadina.\par
\smallskip\noindent\textbf{Variabili} (tutte \emph{Instance Editable}):
\begin{variabili}
ShowBulb & Boolean & sì & Il bulbo visibile: la luce "motivata" (opzionale). \\
\end{variabili}
\noindent In più eredita le variabili di \texttt{BP\_Attivabile} (sezione~\ref{bp:BP_Attivabile}).\par
\smallskip\noindent\textbf{Logica.} \var{Applica(Alpha)}: \nodo{Set Intensity} = 40 $\cdot$ Alpha (parte spenta).\par

\subsection{\texttt{BP\_Ascensore}}\label{bp:BP_Ascensore}
\noindent Rifà \texttt{AActivableMover} (\texttt{Activable.h}). Classe padre: \nodo{BP\_Attivabile}.\par
\smallskip\noindent\textbf{Componenti.} \nodo{Static Mesh} \var{Mesh} (\var{Cube} scala \var{PlatformSize}/100) figlia della root.\par
\noindent\textbf{Origine (pivot).} il \textbf{centro} della piattaforma nella posizione di partenza.\par
\smallskip\noindent\textbf{Variabili} (tutte \emph{Instance Editable}):
\begin{variabili}
EndOffset & Vector & (0, 0, 400) & Punto di arrivo, relativo (si trascina con il rombo nel viewport). \\
PlatformSize & Vector & (300, 300, 30) &  \\
\end{variabili}
\noindent In più eredita le variabili di \texttt{BP\_Attivabile} (sezione~\ref{bp:BP_Attivabile}).\par
\smallskip\noindent\textbf{Logica.} \var{Applica(Alpha)}: \nodo{Set Relative Location} della mesh = \var{EndOffset} $\cdot$ Alpha.\par

\subsection{\texttt{BP\_CresceConDistanza}}\label{bp:BP_CresceConDistanza}
\noindent Rifà \texttt{AScaleByDistanceActor} (\texttt{PlatformerGym.h}). Classe padre: \nodo{Actor}.\par
\smallskip\noindent\textbf{Componenti.} \nodo{Static Mesh} \var{Mesh} (\var{ShapeMesh}, default \var{Cube}), \nodo{Text Render} sopra.\par
\noindent\textbf{Origine (pivot).} il \textbf{centro} della forma.\par
\smallskip\noindent\textbf{Variabili} (tutte \emph{Instance Editable}):
\begin{variabili}
ShapeMesh & StaticMesh (riferimento) &  & La mesh (cubo, sfera...). \\
Material & MaterialInterface (riferimento) &  &  \\
BaseSize & Vector & (300, 300, 30) & Dimensione a scala 1, in cm. \\
MinScale & Float & 0.4 & Scala quando il giocatore è vicino (entro NearDistance). \\
MaxScale & Float & 1.5 & Scala quando il giocatore è lontano (oltre FarDistance). \\
NearDistance & Float & 250 &  \\
FarDistance & Float & 2000 &  \\
ScaleHeight & Boolean & sì & Se falso, cambia solo larghezza e profondità (per le piattaforme: lo spessore resta quello). \\
InterpSpeed & Float & 6 & Quanto in fretta segue la scala voluta (0 = subito). \\
Label & Text &  &  \\
\end{variabili}
\smallskip\noindent\textbf{Logica.} Tick: distanza dal giocatore $\rightarrow$ \nodo{Map Range Clamped} (\var{NearDistance}..\var{FarDistance} $\rightarrow$ \var{MinScale}..\var{MaxScale}) $\rightarrow$ \nodo{FInterp To} (\var{InterpSpeed}) $\rightarrow$ scala = \var{BaseSize}/100 $\cdot$ fattore (se \var{ScaleHeight} è no, la Z resta \var{BaseSize.Z}/100) (Ninja, cap.~7). Nel livello 1 cresce allontanandosi; nel livello 2 le piattaforme \emph{rimpiccioliscono} avvicinandosi: è lo stesso Blueprint con \var{MinScale} $<$ \var{MaxScale}.\par

\subsection{\texttt{BP\_CambiaModello}}\label{bp:BP_CambiaModello}
\noindent Rifà \texttt{AModelCyclerActor} (\texttt{PlatformerGym.h}). Classe padre: \nodo{Actor}.\par
\smallskip\noindent\textbf{Componenti.} \nodo{Static Mesh} \var{Mesh} (scala \var{Size}/100; con \var{Collide} collisione \emph{BlockAll}, altrimenti nessuna), \nodo{Text Render} sopra.\par
\noindent\textbf{Origine (pivot).} il \textbf{centro} della forma.\par
\smallskip\noindent\textbf{Variabili} (tutte \emph{Instance Editable}):
\begin{variabili}
Models & Array di StaticMesh &  & I modelli, in ordine. \\
Materials & Array di MaterialInterface &  & Materiali (uno per modello, oppure uno solo per tutti; vuoto = quelli della mesh). \\
Interval & Float & 1 & Secondi fra un modello e l'altro. \\
Size & Vector & (200, 200, 200) & Dimensione di ogni modello (le forme dell'engine sono di 100 cm). \\
Collide & Boolean & sì & Se vero, ci si può salire sopra (è una piattaforma). \\
Label & Text &  &  \\
\end{variabili}
\smallskip\noindent\textbf{Logica.} \var{Models} = [Cube, Cylinder, Sphere, Cone] (da \texttt{/Engine/BasicShapes}); timer ripetuto ogni \var{Interval}: indice successivo (modulo), \nodo{Set Static Mesh} e \nodo{Set Material} con lo stesso indice in \var{Materials} (Ninja, cap.~7).\par

\subsection{\texttt{BP\_PresenzaCasuale}}\label{bp:BP_PresenzaCasuale}
\noindent Rifà \texttt{ARandomPresenceActor} (\texttt{PlatformerGym.h}). Classe padre: \nodo{Actor}.\par
\smallskip\noindent\textbf{Componenti.} \nodo{Static Mesh} \var{Mesh} (\var{ShapeMesh}, default \var{Cube}; scala \var{Size}/100; \emph{BlockAll}), \nodo{Text Render}.\par
\noindent\textbf{Origine (pivot).} il \textbf{centro} dell'oggetto (per il ponte: centro del ponte, Z = piano $-$ 15).\par
\smallskip\noindent\textbf{Variabili} (tutte \emph{Instance Editable}):
\begin{variabili}
SpawnChance & Float & 0.35 & Probabilità di esserci (0..1). \\
Size & Vector & (1000, 200, 30) & Dimensione in cm. \\
ShapeMesh & StaticMesh (riferimento) &  &  \\
Material & MaterialInterface (riferimento) &  &  \\
Announce & Boolean & sì & Se vero, un messaggio a schermo dice se questa volta c'è. \\
PresentMessage & String & Fortuna! L'oggetto fortunato questa volta c'e' (35\%) &  \\
AbsentMessage & String & L'oggetto fortunato questa volta non c'e'... &  \\
Label & Text &  &  \\
\end{variabili}
\smallskip\noindent\textbf{Logica.} BeginPlay: \nodo{Random Bool with Weight}(\var{SpawnChance}); se esce no, nascosto e senza collisione. Se \var{Announce}, dopo 0,5\,s \var{PresentMessage} o \var{AbsentMessage} (Ninja, cap.~7).\par

\subsection{\texttt{BP\_Mulino}}\label{bp:BP_Mulino}
\noindent Rifà \texttt{AWindmill} (\texttt{Windmill.h}). Classe padre: \nodo{Actor}.\par
\smallskip\noindent\textbf{Componenti.} \var{Torre} (\var{Cylinder} scala $(1{,}6;\ 1{,}6;\ \mathrm{TowerHeight}/100)$ a Z $\mathrm{TowerHeight}/2$), \var{Testa} (Scene a Z \var{TowerHeight}), \var{Mozzo} (Scene a X 110 della testa) con quattro pale (\var{Cube} scala $(0{,}15;\ 0{,}7;\ \mathrm{BladeLength}/100)$, ruotate di 0, 90, 180, 270\textdegree{} in Roll e spostate di $\mathrm{BladeLength}/2+40$ lungo il proprio asse).\par
\noindent\textbf{Origine (pivot).} \textbf{a terra}, al centro della torre.\par
\smallskip\noindent\textbf{Variabili} (tutte \emph{Instance Editable}):
\begin{variabili}
BladeSpeed & Float & 60 & Gradi al secondo delle pale. \\
LookAtPlayer & Boolean & sì & Se vero, la testa si gira verso il player (entro YawLimit). \\
YawLimit & Float & 60 & Quanto può girarsi la testa a destra e a sinistra (gradi, il Clamp). \\
TowerHeight & Float & 900 & Altezza della torre (cm). \\
BladeLength & Float & 450 & Lunghezza di una pala (cm). \\
BladesCollide & Boolean & no & Se vero, le pale hanno collisione (ostacolo). \\
TowerMaterial & MaterialInterface (riferimento) &  &  \\
BladeMaterial & MaterialInterface (riferimento) &  &  \\
\end{variabili}
\smallskip\noindent\textbf{Logica.} pale: \nodo{Add Local Rotation} Roll = \var{BladeSpeed} $\cdot$ dt; se \var{LookAtPlayer}, la testa ruota verso il giocatore con yaw limitato a $\pm$\var{YawLimit} rispetto all'actor (\nodo{Find Look at Rotation}, \nodo{Clamp Angle}) (Ninja, cap.~13).\par

\subsection{\texttt{BP\_CasseFisiche}}\label{bp:BP_CasseFisiche}
\noindent Rifà \texttt{ACubeSpawner} (\texttt{CubeSpawner.h}). Classe padre: \nodo{Actor}.\par
\smallskip\noindent\textbf{Componenti.} una \nodo{Scene} root; le casse sono actor \var{BP\_CassaFisica} (un \var{Cube} con \emph{Simulate Physics}).\par
\noindent\textbf{Origine (pivot).} sopra il centro della griglia di casse (i raggi partono da qui e scendono).\par
\smallskip\noindent\textbf{Variabili} (tutte \emph{Instance Editable}):
\begin{variabili}
Rows & Integer & 3 &  \\
Columns & Integer & 4 &  \\
CubeLayers & Integer & 1 & Quanti cubi uno sopra l'altro in ogni cella. \\
CubeSize & Float & 60 & Lato di un cubo (cm). \\
Spacing & Float & 70 & Distanza fra i centri dei cubi (cm). \\
SimulatePhysics & Boolean & sì & Se vero, i cubi hanno la fisica (si spingono e cadono). \\
CubeMass & Float & 8 & Massa di ogni cubo (kg): leggera = il personaggio li spinge facilmente. \\
Materials & Array di MaterialInterface &  &  \\
DebugTraces & Boolean & no & Disegna in editor le linee dei trace. \\
\end{variabili}
\smallskip\noindent\textbf{Logica.} BeginPlay: doppio \nodo{For Loop} \var{Rows} $\times$ \var{Columns}; per ogni cella un \nodo{Line Trace} verso il basso trova il terreno, poi \var{CubeLayers} casse impilate di lato \var{CubeSize}, passo \var{Spacing}, massa \var{CubeMass}, materiale alternato da \var{Materials} (la lezione dei due cicli for).\par

\subsection{\texttt{BP\_Staccionata}}\label{bp:BP_Staccionata}
\noindent Rifà \texttt{ASplineFence} (\texttt{LevelDressing.h}). Classe padre: \nodo{Actor}.\par
\smallskip\noindent\textbf{Componenti.} \nodo{Spline} \var{Percorso}, \nodo{Instanced Static Mesh} \var{Pali} (\var{PostMesh}, default \var{Cube} scala \var{PostSize}).\par
\noindent\textbf{Origine (pivot).} il primo punto della spline.\par
\smallskip\noindent\textbf{Variabili} (tutte \emph{Instance Editable}):
\begin{variabili}
PostMesh & StaticMesh (riferimento) &  & Mesh dei paletti (vuota = cilindro). \\
Spacing & Float & 200 & Un paletto ogni Spacing cm lungo la spline. \\
PostSize & Vector & (14, 14, 110) & Dimensioni di un paletto (cm). \\
Rails & Integer & 2 & Quante traverse fra un paletto e l'altro (0 = nessuna). \\
RailThickness & Float & 8 & Spessore delle traverse (cm). \\
Material & MaterialInterface (riferimento) &  &  \\
Collide & Boolean & sì & Se vero, la staccionata blocca il personaggio. \\
\end{variabili}
\smallskip\noindent\textbf{Logica.} Construction Script: un palo ogni \var{Spacing} cm lungo la spline, \var{Rails} traverse fra i pali con spessore \var{RailThickness} (Ninja, cap.~13). I punti della spline sono nelle tabelle (coordinate \emph{locali}).\par

\subsection{\texttt{BP\_Spargi}}\label{bp:BP_Spargi}
\noindent Rifà \texttt{AScatterArea} (\texttt{LevelDressing.h}). Classe padre: \nodo{Actor}.\par
\smallskip\noindent\textbf{Componenti.} \nodo{Box Collision} \var{Area} (extent \var{AreaExtent}, solo per vederla nell'editor), \nodo{Instanced Static Mesh} \var{Istanze} (\var{Mesh}).\par
\noindent\textbf{Origine (pivot).} il \textbf{centro} dell'area (Z: sopra il terreno, i raggi scendono).\par
\smallskip\noindent\textbf{Variabili} (tutte \emph{Instance Editable}):
\begin{variabili}
Mesh & StaticMesh (riferimento) &  &  \\
AreaExtent & Vector & (1000, 1000, 500) & Metà delle dimensioni dell'area (cm). \\
Count & Integer & 30 &  \\
Seed & Integer & 7 & Seme del caso: stesso seme = stessa disposizione. \\
MinScale & Float & 0.8 &  \\
MaxScale & Float & 1.3 &  \\
MinDistance & Float & 150 & Distanza minima fra due istanze (cm), come il raggio del foliage. \\
AlignToGround & Boolean & sì & Se vero, ogni istanza è appoggiata al terreno con un trace verso il basso. \\
Collide & Boolean & no & Se vero, le istanze hanno collisione (gli alberi sì, l'erba no). \\
CullDistance & Float & 0 & Oltre questa distanza le istanze spariscono (0 = mai). \\
\end{variabili}
\smallskip\noindent\textbf{Logica.} Construction Script con un \nodo{Random Stream} di seme \var{Seed}: \var{Count} punti a caso nell'area, scartando quelli a meno di \var{MinDistance} da un altro, \nodo{Line Trace} verso il basso (se \var{AlignToGround}), scala fra \var{MinScale} e \var{MaxScale}, yaw a caso, \nodo{Add Instance} (Ninja, cap.~13). Negli alberi: mesh \var{SM\_Chestnut\_4}.\par

\subsection{\texttt{BP\_LuceTremolante}}\label{bp:BP_LuceTremolante}
\noindent Rifà \texttt{AFlickeringLight} (\texttt{LevelDressing.h}). Classe padre: \nodo{Actor}.\par
\smallskip\noindent\textbf{Componenti.} \nodo{Point Light} (root o figlia; Movable, temperatura \var{Temperature}, raggio \var{AttenuationRadius}, \emph{Source Radius} \var{SourceRadius}).\par
\noindent\textbf{Origine (pivot).} dove sta la fiamma (nei bracieri: 150\,cm sopra il piano).\par
\smallskip\noindent\textbf{Variabili} (tutte \emph{Instance Editable}):
\begin{variabili}
BaseIntensity & Float & 60 & Intensità media (candele). \\
FlickerAmount & Float & 0.25 & Quanto varia l'intensità (0 = ferma, 0.5 = ±50\%). \\
FlickerSpeed & Float & 6 & Velocità del tremolio. \\
PositionJitter & Float & 4 & Quanto si sposta la sorgente (cm): le ombre "respirano" come con una fiamma vera. \\
Temperature & Float & 1900 & Temperatura colore (K): 1900 fuoco, 2700 lampadina calda, 6500 luce fredda. \\
AttenuationRadius & Float & 900 &  \\
SourceRadius & Float & 15 & Raggio della sorgente (cm): più grande = ombre più morbide. \\
CastShadows & Boolean & sì &  \\
GlowActors & Array di Actor &  & Attori con un materiale emissivo da far tremolare insieme alla luce (braci, fiamme). \\
GlowParameter & Name & Intensita & Nome del parametro scalare del materiale emissivo da animare. \\
GlowBaseValue & Float & 8 & Valore del parametro quando la luce è all'intensità media. \\
\end{variabili}
\smallskip\noindent\textbf{Logica.} Tick: intensità = \var{BaseIntensity} $\cdot$ (1 + rumore $\cdot$ \var{FlickerAmount}) con \nodo{FInterp To} verso un bersaglio casuale (\var{FlickerSpeed}); piccolo spostamento casuale di \var{PositionJitter} cm; sugli actor in \var{GlowActors} imposta il parametro \var{GlowParameter} del materiale emissivo in proporzione (Ninja, cap.~13). Nel braciere \var{GlowActors} = la sfera delle braci.\par



% =====================================================================
\addtocontents{toc}{\protect\setcounter{tocdepth}{2}}
\section{Livello 1 -- ``Si impara''}

Mattina, un'isola sull'acqua. Si impara a muoversi e si raccolgono le prime due abilità: il doppio salto (serve per la parete di 3,3\,m) e il generatore di casse (serve per la parete di 4,6\,m, che nemmeno il doppio salto supera). Il ponte fortunato c'è solo nel 35\% delle partite: altrimenti si passa dai sassi. Cadere in acqua costa una vita.

Livello: \var{L\_GD13\_Livello1\_BP}. Traguardo $\rightarrow$ \var{L\_GD13\_Livello2\_BP}.

%% generato da strumenti/estrai_livelli.py: non modificare a mano
\subsection{Luce, atmosfera e partenza}
\begin{itemize}
\item \nodo{Directional Light} (il sole): Mobility \emph{Movable}, Rotation pitch -32, yaw -35; Intensity 9 lux; \emph{Use Temperature} sì, Temperature 5200 K; \emph{Atmosphere Sun Light} sì; Source Angle 1.
\item \nodo{Sky Atmosphere} (valori di default).
\item \nodo{Sky Light}: Movable, \emph{Real Time Capture} sì, Intensity 1.
\item \nodo{Exponential Height Fog} a Z $-200$: Fog Density 0,01.
\item \nodo{Post Process Volume}: \emph{Infinite Extent (Unbound)} sì; spuntate e impostate Auto Exposure Bias 0,3, Bloom Intensity 0,5, Vignette Intensity 0,25.
\item \nodo{Player Start} in (-500; 0; 120), rotazione 0.
\end{itemize}

\subsection{Partenza e primo tratto}
\begin{oggetti}
\nota{acqua tutto intorno: chi cade perde una vita e torna all'ultimo posto sicuro (la zona sta SOPRA l'acqua, che ha collisione: si "affoga" prima di toccarla)}
1 & Cubo & Acqua & (4500; 0; -310) & (17000; 9000; 20) & MI\_\allowbreak{}AIV\_\allowbreak{}Acqua \\
2 & BP\_\allowbreak{}ZonaCaduta & Zona\_\allowbreak{}Caduta & (4500; 0; -200) & -- & ZoneExtent = (8500; 4500; 150); ShowVisual = no \\
\end{oggetti}

\subsection{Isola di partenza}
\begin{oggetti}
3 & Pedana & Isola\_\allowbreak{}Partenza & (0; 0; -150) & (1600; 1600; 300) & MI\_\allowbreak{}AIV\_\allowbreak{}Erba, si cammina a Z = 0 \\
4 & Text Render & Scritta & (-150; 0; 330) & yaw 180 & ``LIVELLO 1 / salta, raccogli, scopri'', World Size 55, colore (1; 1; 1) \\
5 & Text Render & Scritta & (-150; 0; 230) & yaw 180 & ``3 VITE: se cadi in acqua torni all'ultimo posto sicuro / (a zero vite si ricomincia dal livello 1)'', World Size 26, colore (1; 0,55; 0,65) \\
\nota{il faro: piu' ti allontani, piu' diventa grande (cosi' lo vedi sempre)}
6 & Cubo & Faro\_\allowbreak{}Palo & (-650; 650; 200) & (30; 30; 400) & MI\_\allowbreak{}AIV\_\allowbreak{}LegnoScuro \\
7 & BP\_\allowbreak{}CresceConDistanza & Faro\_\allowbreak{}CresceConLaDistanza & (-650; 650; 450) & -- & ShapeMesh = Sphere; Material = MI\_\allowbreak{}AIV\_\allowbreak{}Abilita; BaseSize = (100; 100; 100); MinScale = 1; MaxScale = 7; NearDistance = 400; FarDistance = 9000; ScaleHeight = sì; Label = ``FARO: piu' sei lontano, piu' cresce'' \\
8 & BP\_\allowbreak{}Moneta & Moneta\_\allowbreak{}Isola\_\allowbreak{}1 & (-300; -300; 90) & -- & Value = 1; Material = MI\_\allowbreak{}AIV\_\allowbreak{}Moneta; GlowIntensity = 4; GlowColor = (1; 0,75; 0,3) \\
9 & BP\_\allowbreak{}Moneta & Moneta\_\allowbreak{}Isola\_\allowbreak{}2 & (-50; -300; 90) & -- & Value = 1; Material = MI\_\allowbreak{}AIV\_\allowbreak{}Moneta; GlowIntensity = 4; GlowColor = (1; 0,75; 0,3) \\
10 & BP\_\allowbreak{}Moneta & Moneta\_\allowbreak{}Isola\_\allowbreak{}3 & (200; -300; 90) & -- & Value = 1; Material = MI\_\allowbreak{}AIV\_\allowbreak{}Moneta; GlowIntensity = 4; GlowColor = (1; 0,75; 0,3) \\
\nota{pozzanghera: il materiale fisico dell'acqua cambia il suono dei passi}
11 & Cubo & Pozzanghera & (-250; 350; 1) & (360; 260; 2) & MI\_\allowbreak{}AIV\_\allowbreak{}Acqua \\
12 & Decal & Pozza\_\allowbreak{}Bordo & (-250; 350; 10) & rot (-90; 0; 0), Decal Size (30; 140; 190) & MI\_\allowbreak{}AIV\_\allowbreak{}Pozza \\
\nota{mulino: le pale ruotano, la testa ti segue (con un limite all'angolo)}
13 & BP\_\allowbreak{}Mulino & Mulino & (-450; -650; 0) & yaw 60 & BladeSpeed = 50; LookAtPlayer = sì; YawLimit = 60; TowerHeight = 900; BladeLength = 420; BladesCollide = no; TowerMaterial = MI\_\allowbreak{}AIV\_\allowbreak{}IntonacoChiaro; BladeMaterial = MI\_\allowbreak{}AIV\_\allowbreak{}LegnoChiaro \\
\nota{casse fisiche: spawn con due cicli for e un trace verso il basso}
14 & BP\_\allowbreak{}CasseFisiche & Casse\_\allowbreak{}Fisiche & (250; 550; 300) & -- & Rows = 3; Columns = 3; CubeLayers = 2; CubeSize = 60; Spacing = 66; SimulatePhysics = sì; CubeMass = 6; Materials = [MI\_\allowbreak{}AIV\_\allowbreak{}LegnoChiaro, MI\_\allowbreak{}AIV\_\allowbreak{}LegnoScuro] \\
\nota{staccionata lungo una spline (un paletto ogni 200 cm)}
15 & BP\_\allowbreak{}Staccionata & Staccionata & (-780; -780; 0) & -- & Spacing = 200; Rails = 2; Material = MI\_\allowbreak{}AIV\_\allowbreak{}LegnoScuro; Collide = sì; punti della spline (locali): (0; 0; 0); (0; 1560; 0); (1560; 1560; 0) \\
16 & Static Mesh & Cespuglio\_\allowbreak{}0 & (-600; 620; 0) & scala (1,3; 1,3; 1,3) & SM\_\allowbreak{}Groundplant\_\allowbreak{}12; Collision: No Collision \\
17 & Static Mesh & Cespuglio\_\allowbreak{}1 & (500; -700; 0) & scala (1,3; 1,3; 1,3) yaw 70 & SM\_\allowbreak{}Groundplant\_\allowbreak{}12; Collision: No Collision \\
18 & Static Mesh & Cespuglio\_\allowbreak{}2 & (-700; 100; 0) & scala (1,3; 1,3; 1,3) yaw 140 & SM\_\allowbreak{}Groundplant\_\allowbreak{}12; Collision: No Collision \\
\nota{alberi sulle isolette intorno (foliage "a mano")}
19 & Pedana & Isoletta\_\allowbreak{}Nord & (1500; 2600; -150) & (2400; 1400; 300) & MI\_\allowbreak{}AIV\_\allowbreak{}Erba, si cammina a Z = 0 \\
20 & Pedana & Isoletta\_\allowbreak{}Sud & (4500; -2900; -150) & (3200; 1600; 300) & MI\_\allowbreak{}AIV\_\allowbreak{}Erba, si cammina a Z = 0 \\
21 & BP\_\allowbreak{}Spargi & Alberi\_\allowbreak{}Nord & (1500; 2600; 300) & -- & Mesh = SM\_\allowbreak{}Chestnut\_\allowbreak{}4; AreaExtent = (1100; 600; 600); Count = 5; Seed = 11; MinScale = 0,55; MaxScale = 0,8; MinDistance = 600; AlignToGround = sì; Collide = sì; CullDistance = 0 \\
22 & BP\_\allowbreak{}Spargi & Alberi\_\allowbreak{}Sud & (4500; -2900; 300) & -- & Mesh = SM\_\allowbreak{}Chestnut\_\allowbreak{}4; AreaExtent = (1500; 700; 600); Count = 7; Seed = 23; MinScale = 0,55; MaxScale = 0,8; MinDistance = 600; AlignToGround = sì; Collide = sì; CullDistance = 0 \\
\end{oggetti}

\subsection{Cancello: pedana a pressione $\rightarrow$ porta doppia + luce}
\begin{oggetti}
\nota{muro con un'apertura al centro}
23 & Cubo & Cancello\_\allowbreak{}Muro\_\allowbreak{}S & (780; -460; 200) & (60; 680; 400) & MI\_\allowbreak{}AIV\_\allowbreak{}Pietra \\
24 & Cubo & Cancello\_\allowbreak{}Muro\_\allowbreak{}N & (780; 460; 200) & (60; 680; 400) & MI\_\allowbreak{}AIV\_\allowbreak{}Pietra \\
25 & Cubo & Cancello\_\allowbreak{}Architrave & (780; 0; 350) & (60; 240; 100) & MI\_\allowbreak{}AIV\_\allowbreak{}Pietra \\
26 & BP\_\allowbreak{}Porta & Porta\_\allowbreak{}Doppia & (780; 0; 0) & -- & DoorWidth = 240; DoorHeight = 300; DoorThickness = 16; Sliding = no; OpenAngle = 100; TransitionTime = 1,2; DoorMaterial = MI\_\allowbreak{}AIV\_\allowbreak{}LegnoDipinto \\
27 & BP\_\allowbreak{}PiastraPressione & Pedana\_\allowbreak{}Pressione & (450; -300; 0) & -- & StayPressed = sì; OnMaterial = MI\_\allowbreak{}AIV\_\allowbreak{}Acceso; OffMaterial = MI\_\allowbreak{}AIV\_\allowbreak{}Spento; Targets = [Porta\_\allowbreak{}Doppia] \\
28 & BP\_\allowbreak{}LuceAttivabile & Lampada\_\allowbreak{}Cancello & (720; 0; 360) & -- & TransitionTime = 0,6; ListenTo = [Pedana\_\allowbreak{}Pressione] \\
29 & Text Render & Scritta & (450; -300; 140) & yaw 180 & ``sali sulla pedana'', World Size 35, colore (1; 1; 1) \\
\end{oggetti}

\subsection{Gradini}
\begin{oggetti}
30 & Pedana & P1 & (1100; 0; -20) & (400; 400; 40) & MI\_\allowbreak{}AIV\_\allowbreak{}Percorso, si cammina a Z = 0 \\
31 & Pedana & P2\_\allowbreak{}Gradino & (1550; 0; 40) & (300; 300; 40) & MI\_\allowbreak{}AIV\_\allowbreak{}PercorsoScuro, si cammina a Z = 60 \\
32 & Pedana & P3\_\allowbreak{}Gradino & (1900; 0; 100) & (300; 300; 40) & MI\_\allowbreak{}AIV\_\allowbreak{}Percorso, si cammina a Z = 120 \\
33 & Pedana & P4\_\allowbreak{}Gradino & (2300; 0; 160) & (300; 300; 40) & MI\_\allowbreak{}AIV\_\allowbreak{}PercorsoScuro, si cammina a Z = 180 \\
34 & BP\_\allowbreak{}Moneta & Moneta\_\allowbreak{}P2 & (1550; 0; 150) & -- & Value = 1; Material = MI\_\allowbreak{}AIV\_\allowbreak{}Moneta; GlowIntensity = 4; GlowColor = (1; 0,75; 0,3) \\
35 & BP\_\allowbreak{}Moneta & Moneta\_\allowbreak{}P3 & (1900; 0; 210) & -- & Value = 1; Material = MI\_\allowbreak{}AIV\_\allowbreak{}Moneta; GlowIntensity = 4; GlowColor = (1; 0,75; 0,3) \\
36 & BP\_\allowbreak{}Moneta & Moneta\_\allowbreak{}P4 & (2300; 0; 270) & -- & Value = 1; Material = MI\_\allowbreak{}AIV\_\allowbreak{}Moneta; GlowIntensity = 4; GlowColor = (1; 0,75; 0,3) \\
\end{oggetti}

\subsection{Piattaforma mobile (suono con pitch diverso all'andata e al ritorno)}
\begin{oggetti}
37 & BP\_\allowbreak{}PiattaformaMobile & Mobile\_\allowbreak{}1 & (2650; 0; 165) & -- & EndOffset = (600; 0; 0); TravelTime = 2,5; PauseTime = 0,8; StartDelay = 0; PlatformSize = (250; 250; 30); PlatformMaterial = MI\_\allowbreak{}AIV\_\allowbreak{}Salto; MoveSound = SfxPlatformMoving\_\allowbreak{}Cue \\
38 & Text Render & Scritta & (2450; 0; 420) & yaw 180 & ``aspetta la piattaforma'', World Size 35, colore (1; 1; 1) \\
\end{oggetti}

\subsection{Corridoio con la zona velocità e salto lungo}
\begin{oggetti}
39 & Pedana & P5\_\allowbreak{}Corridoio & (3900; 0; 160) & (900; 400; 40) & MI\_\allowbreak{}AIV\_\allowbreak{}Percorso, si cammina a Z = 180 \\
40 & BP\_\allowbreak{}ZonaModificatore & Zona\_\allowbreak{}Velocita & (3650; 0; 280) & -- & ZoneExtent = (180; 180; 100); WalkSpeedMultiplier = 1,8; DurationAfterExit = 4; Label = ``VELOCITA x1.8''; VisualMaterial = MI\_\allowbreak{}AIV\_\allowbreak{}Salto \\
41 & Text Render & Scritta & (4250; 0; 420) & yaw 180 & ``prendi la rincorsa!'', World Size 40, colore (1; 1; 1) \\
\end{oggetti}

\subsection{IL DOPPIO SALTO}
\begin{oggetti}
42 & Pedana & P6\_\allowbreak{}DoppioSalto & (5000; 0; 160) & (500; 500; 40) & MI\_\allowbreak{}AIV\_\allowbreak{}Salto, si cammina a Z = 180 \\
43 & BP\_\allowbreak{}PickupAbilita & Oggetto\_\allowbreak{}DoppioSalto & (5000; 0; 290) & -- & Ability = DoppioSalto; Material = MI\_\allowbreak{}AIV\_\allowbreak{}Abilita; GlowIntensity = 20; GlowColor = (0,3; 0,7; 1); PickupRadius = 80; Label = ``DOPPIO SALTO''; Message = ``Doppio salto! Premi salto di nuovo in aria: puoi anche cambiare direzione'' \\
44 & Spot Light & Luce\_\allowbreak{}DoppioSalto & (5000; 0; 900) & pitch -90 yaw 0 & 400 cd, raggio 1200, colore (0,5; 0,75; 1), coni 12/22° \\
\nota{parete di 3,3 m: con un salto solo (2,15 m) non ci si arriva}
45 & Pedana & P7\_\allowbreak{}Parete & (5550; 0; 345) & (500; 500; 330) & MI\_\allowbreak{}AIV\_\allowbreak{}PercorsoScuro, si cammina a Z = 510 \\
46 & Text Render & Scritta & (5250; 0; 560) & yaw 180 & ``troppo alta? salta due volte'', World Size 35, colore (1; 1; 1) \\
\nota{salti in cui, se sbagli, ti correggi in aria con il secondo salto}
47 & Pedana & P8 & (6250; 0; 490) & (300; 300; 40) & MI\_\allowbreak{}AIV\_\allowbreak{}Percorso, si cammina a Z = 510 \\
48 & Pedana & P9 & (6900; 200; 580) & (300; 300; 40) & MI\_\allowbreak{}AIV\_\allowbreak{}PercorsoScuro, si cammina a Z = 600 \\
49 & Pedana & P10 & (7550; -100; 490) & (300; 300; 40) & MI\_\allowbreak{}AIV\_\allowbreak{}Percorso, si cammina a Z = 510 \\
50 & BP\_\allowbreak{}Moneta & Moneta\_\allowbreak{}P8 & (6250; 0; 600) & -- & Value = 1; Material = MI\_\allowbreak{}AIV\_\allowbreak{}Moneta; GlowIntensity = 4; GlowColor = (1; 0,75; 0,3) \\
51 & BP\_\allowbreak{}Moneta & Moneta\_\allowbreak{}P9 & (6900; 200; 690) & -- & Value = 1; Material = MI\_\allowbreak{}AIV\_\allowbreak{}Moneta; GlowIntensity = 4; GlowColor = (1; 0,75; 0,3) \\
52 & BP\_\allowbreak{}Moneta & Moneta\_\allowbreak{}P10 & (7550; -100; 600) & -- & Value = 1; Material = MI\_\allowbreak{}AIV\_\allowbreak{}Moneta; GlowIntensity = 4; GlowColor = (1; 0,75; 0,3) \\
53 & Text Render & Scritta & (6250; 0; 780) & yaw 180 & ``se manchi il salto, risalta in aria / e correggi la direzione'', World Size 32, colore (1; 1; 1) \\
\end{oggetti}

\subsection{GENERATORE DI CASSE: la parete di 4,6 m}
\begin{oggetti}
54 & Pedana & P11\_\allowbreak{}Casse & (8300; 0; 490) & (700; 600; 40) & MI\_\allowbreak{}AIV\_\allowbreak{}Salto, si cammina a Z = 510 \\
55 & BP\_\allowbreak{}PickupAbilita & Oggetto\_\allowbreak{}Casse & (8100; 0; 620) & -- & Ability = GeneratoreCasse; Material = MI\_\allowbreak{}AIV\_\allowbreak{}Abilita; GlowIntensity = 20; GlowColor = (0,3; 0,7; 1); PickupRadius = 80; Label = ``GENERATORE DI CASSE''; Message = ``Generatore di casse! Premi F: fino a 3 casse, poi 5 secondi di ricarica'' \\
56 & Spot Light & Luce\_\allowbreak{}Oggetto\_\allowbreak{}Casse & (8100; 0; 1210) & pitch -90 yaw 0 & 300 cd, raggio 1000, colore (0,5; 0,75; 1), coni 12/22° \\
57 & Pedana & P12\_\allowbreak{}ParetePiuAlta & (8900; 0; 590) & (500; 600; 760) & MI\_\allowbreak{}AIV\_\allowbreak{}PercorsoScuro, si cammina a Z = 970 \\
58 & Text Render & Scritta & (8300; 0; 760) & yaw 180 & ``parete di 4,6 m: nemmeno il doppio salto basta / F davanti al muro: 2 casse una sull'altra, salici e salta due volte'', World Size 32, colore (1; 0,9; 0,5) \\
\end{oggetti}

\subsection{PONTE FORTUNATO (c'e' solo il 35\% delle volte) e la strada lunga}
\begin{oggetti}
59 & BP\_\allowbreak{}PresenzaCasuale & Ponte\_\allowbreak{}Fortunato & (9625; 0; 955) & -- & SpawnChance = 0,35; Size = (950; 250; 30); Material = MI\_\allowbreak{}AIV\_\allowbreak{}Moneta; Label = ``PONTE FORTUNATO (35\%)'' \\
60 & Pedana & Sasso\_\allowbreak{}1 & (9450; -450; 950) & (200; 200; 40) & MI\_\allowbreak{}AIV\_\allowbreak{}PercorsoScuro, si cammina a Z = 970 \\
61 & Pedana & Sasso\_\allowbreak{}2 & (9850; -450; 980) & (200; 200; 40) & MI\_\allowbreak{}AIV\_\allowbreak{}PercorsoScuro, si cammina a Z = 1000 \\
62 & Text Render & Scritta & (9150; 0; 1230) & yaw 180 & ``il ponte c'e' solo 35 volte su 100 / se non c'e': i sassi a sinistra'', World Size 32, colore (1; 0,9; 0,5) \\
\end{oggetti}

\subsection{Il totem che cambia modello ogni secondo}
\begin{oggetti}
63 & Pedana & P13 & (10400; 0; 950) & (600; 600; 40) & MI\_\allowbreak{}AIV\_\allowbreak{}Percorso, si cammina a Z = 970 \\
64 & BP\_\allowbreak{}CambiaModello & Totem\_\allowbreak{}CambiaModello & (10450; 220; 1040) & -- & Interval = 1; Size = (120; 120; 140); Collide = sì; Materials = [MI\_\allowbreak{}AIV\_\allowbreak{}Salto, MI\_\allowbreak{}AIV\_\allowbreak{}Abilita, MI\_\allowbreak{}AIV\_\allowbreak{}Moneta, MI\_\allowbreak{}AIV\_\allowbreak{}Pericolo]; Label = ``CAMBIA MODELLO OGNI SECONDO'' \\
65 & BP\_\allowbreak{}Moneta & Moneta\_\allowbreak{}P13 & (10350; -150; 1060) & -- & Value = 2; Material = MI\_\allowbreak{}AIV\_\allowbreak{}Moneta; GlowIntensity = 4; GlowColor = (1; 0,75; 0,3) \\
\end{oggetti}

\subsection{Traguardo}
\begin{oggetti}
66 & Pedana & P14\_\allowbreak{}Arrivo & (11150; 0; 950) & (600; 600; 40) & MI\_\allowbreak{}AIV\_\allowbreak{}Traguardo, si cammina a Z = 970 \\
67 & BP\_\allowbreak{}Traguardo & Traguardo & (11200; 0; 1070) & -- & ZoneExtent = (150; 150; 100); NextLevel = L\_\allowbreak{}GD13\_\allowbreak{}Livello2\_\allowbreak{}BP; Message = ``Livello 1 completato! Le monete sono al sicuro''; Delay = 1,2; ResetDeathCount = no; Label = ``TRAGUARDO''; VisualMaterial = MI\_\allowbreak{}AIV\_\allowbreak{}Traguardo \\
68 & Point Light & Luce\_\allowbreak{}Traguardo & (11200; 0; 1230) & -- & 30 cd, raggio 700, colore (0,4; 1; 0,5), niente ombre \\
\end{oggetti}



\textbf{Collaudo del livello 1.} (1) La pedana a pressione apre la porta e accende la lampada. (2) Senza doppio salto P7 non si raggiunge; dopo averlo preso, sì. (3) Davanti a P12 bastano due casse una sull'altra e un doppio salto. (4) Riavviate qualche volta: il ponte a volte c'è e a volte no, e i sassi portano comunque avanti. (5) Cadete in acqua tre volte: game over e si ricomincia da \var{L\_GD13\_Livello1\_BP}.

% =====================================================================
\section{Livello 2 -- ``Ostacoli e congegni''}

Pomeriggio. Trampolino, torretta con i ripari, piattaforme che crollano, leva e ascensore, la sezione del wall run fra due muri blu (con un pavimento e una pedana per tornare su se si cade), lava da 10\,HP al secondo, la sbarra che spinge, il ghiaccio o il ponte a pagamento, il jetpack per il salto lunghissimo, le piattaforme che rimpiccioliscono e la torretta i cui proiettili fanno ripartire il livello.

Livello: \var{L\_GD13\_Livello2\_BP}. Traguardo $\rightarrow$ \var{L\_GD13\_Livello3\_BP}.

%% generato da strumenti/estrai_livelli.py: non modificare a mano
\subsection{Luce, atmosfera e partenza}
\begin{itemize}
\item \nodo{Directional Light} (il sole): Mobility \emph{Movable}, Rotation pitch -38, yaw 40; Intensity 8 lux; \emph{Use Temperature} sì, Temperature 5600 K; \emph{Atmosphere Sun Light} sì; Source Angle 1.
\item \nodo{Sky Atmosphere} (valori di default).
\item \nodo{Sky Light}: Movable, \emph{Real Time Capture} sì, Intensity 1.
\item \nodo{Exponential Height Fog} a Z $-200$: Fog Density 0,02.
\item \nodo{Post Process Volume}: \emph{Infinite Extent (Unbound)} sì; spuntate e impostate Auto Exposure Bias 0,2, Bloom Intensity 0,6, Vignette Intensity 0,3.
\item \nodo{Player Start} in (-150; 0; 120), rotazione 0.
\end{itemize}

\subsection{Partenza e primo tratto}
\begin{oggetti}
1 & Text Render & Scritta & (150; 0; 330) & yaw 180 & ``LIVELLO 2 / ostacoli e congegni'', World Size 55, colore (1; 1; 1) \\
2 & BP\_\allowbreak{}ZonaCaduta & Zona\_\allowbreak{}Caduta & (9500; 0; -600) & -- & ZoneExtent = (12500; 4000; 150); ShowVisual = no \\
3 & Pedana & P0\_\allowbreak{}Partenza & (0; 0; -20) & (800; 800; 40) & MI\_\allowbreak{}AIV\_\allowbreak{}Percorso, si cammina a Z = 0 \\
\end{oggetti}

\subsection{Pedana di lancio (Launch Character)}
\begin{oggetti}
4 & Pedana & P1\_\allowbreak{}Lancio & (800; 0; -20) & (400; 400; 40) & MI\_\allowbreak{}AIV\_\allowbreak{}PercorsoScuro, si cammina a Z = 0 \\
5 & BP\_\allowbreak{}Trampolino & Pedana\_\allowbreak{}Lancio & (800; 0; 10) & -- & LaunchVelocity = (350; 0; 1300); OverrideHorizontal = sì; PadMaterial = MI\_\allowbreak{}AIV\_\allowbreak{}Salto \\
6 & Pedana & P2\_\allowbreak{}Alto & (1500; 0; 250) & (500; 500; 500) & MI\_\allowbreak{}AIV\_\allowbreak{}Percorso, si cammina a Z = 500 \\
7 & BP\_\allowbreak{}Moneta & Moneta\_\allowbreak{}P2 & (1500; 0; 590) & -- & Value = 1; Material = MI\_\allowbreak{}AIV\_\allowbreak{}Moneta; GlowIntensity = 4; GlowColor = (1; 0,75; 0,3) \\
\end{oggetti}

\subsection{La torretta (con ripari e una cura)}
\begin{oggetti}
8 & Pedana & P3\_\allowbreak{}Tiro & (2500; 0; 480) & (1200; 600; 40) & MI\_\allowbreak{}AIV\_\allowbreak{}Percorso, si cammina a Z = 500 \\
9 & Cubo & Riparo\_\allowbreak{}1 & (2150; -80; 600) & (80; 300; 200) & MI\_\allowbreak{}AIV\_\allowbreak{}LegnoScuro \\
10 & Cubo & Riparo\_\allowbreak{}2 & (2600; 120; 600) & (80; 300; 200) & MI\_\allowbreak{}AIV\_\allowbreak{}LegnoScuro \\
11 & Pedana & Torretta\_\allowbreak{}Base & (2900; -650; 250) & (300; 300; 700) & MI\_\allowbreak{}AIV\_\allowbreak{}Pericolo, si cammina a Z = 600 \\
12 & BP\_\allowbreak{}Torretta & Torretta & (2900; -650; 600) & yaw 90 & Range = 1800; FireInterval = 1,6; TurnSpeed = 3; ProjectileSpeed = 1400; ProjectileDamage = 20; LeadTarget = no; BodyMaterial = MI\_\allowbreak{}AIV\_\allowbreak{}Torretta; ProjectileMaterial = MI\_\allowbreak{}AIV\_\allowbreak{}Proiettile \\
13 & BP\_\allowbreak{}Moneta & Moneta\_\allowbreak{}P3\_\allowbreak{}1 & (2000; 150; 590) & -- & Value = 1; Material = MI\_\allowbreak{}AIV\_\allowbreak{}Moneta; GlowIntensity = 4; GlowColor = (1; 0,75; 0,3) \\
14 & BP\_\allowbreak{}Moneta & Moneta\_\allowbreak{}P3\_\allowbreak{}2 & (2400; 150; 590) & -- & Value = 1; Material = MI\_\allowbreak{}AIV\_\allowbreak{}Moneta; GlowIntensity = 4; GlowColor = (1; 0,75; 0,3) \\
15 & BP\_\allowbreak{}PickupSalute & Cura\_\allowbreak{}P3 & (3000; 200; 590) & -- & HealAmount = 40; Material = MI\_\allowbreak{}AIV\_\allowbreak{}Vita; GlowColor = (0,3; 1; 0,4); Label = ``+VITA'' \\
16 & Text Render & Scritta & (2000; 0; 800) & yaw 180 & ``la torretta ti vede: usa i ripari'', World Size 35, colore (1; 1; 1) \\
\end{oggetti}

\subsection{Piattaforme che crollano}
\begin{oggetti}
17 & BP\_\allowbreak{}PiattaformaCrolla & Crolla\_\allowbreak{}1 & (3400; 0; 485) & -- & PlatformSize = (250; 250; 30); ShakeTime = 0,7; RespawnTime = 3; PlatformMaterial = MI\_\allowbreak{}AIV\_\allowbreak{}Salto; WarningMaterial = MI\_\allowbreak{}AIV\_\allowbreak{}Pericolo; CrumbleSound = Rock\_\allowbreak{}Cue \\
18 & BP\_\allowbreak{}PiattaformaCrolla & Crolla\_\allowbreak{}2 & (3800; 0; 485) & -- & PlatformSize = (250; 250; 30); ShakeTime = 0,7; RespawnTime = 3; PlatformMaterial = MI\_\allowbreak{}AIV\_\allowbreak{}Salto; WarningMaterial = MI\_\allowbreak{}AIV\_\allowbreak{}Pericolo; CrumbleSound = Rock\_\allowbreak{}Cue \\
19 & BP\_\allowbreak{}PiattaformaCrolla & Crolla\_\allowbreak{}3 & (4200; 0; 485) & -- & PlatformSize = (250; 250; 30); ShakeTime = 0,7; RespawnTime = 3; PlatformMaterial = MI\_\allowbreak{}AIV\_\allowbreak{}Salto; WarningMaterial = MI\_\allowbreak{}AIV\_\allowbreak{}Pericolo; CrumbleSound = Rock\_\allowbreak{}Cue \\
20 & BP\_\allowbreak{}Moneta & Moneta\_\allowbreak{}Crolla & (3800; 0; 590) & -- & Value = 1; Material = MI\_\allowbreak{}AIV\_\allowbreak{}Moneta; GlowIntensity = 4; GlowColor = (1; 0,75; 0,3) \\
\end{oggetti}

\subsection{Leva ("premi E") e ascensore}
\begin{oggetti}
21 & Pedana & P4\_\allowbreak{}Leva & (4750; 0; 480) & (600; 600; 40) & MI\_\allowbreak{}AIV\_\allowbreak{}PercorsoScuro, si cammina a Z = 500 \\
22 & BP\_\allowbreak{}Ascensore & Ascensore & (5250; 0; 485) & -- & EndOffset = (0; 0; 450); PlatformSize = (300; 300; 30); TransitionTime = 2,5; Material (mesh) = MI\_\allowbreak{}AIV\_\allowbreak{}Salto \\
23 & BP\_\allowbreak{}Leva & Leva\_\allowbreak{}Ascensore & (4900; -220; 500) & yaw 180 & Prompt = ``Chiama l'ascensore''; Toggle = sì; Label = ``LEVA''; OnMaterial = MI\_\allowbreak{}AIV\_\allowbreak{}Acceso; OffMaterial = MI\_\allowbreak{}AIV\_\allowbreak{}Spento; Targets = [Ascensore] \\
\end{oggetti}

\subsection{WALL RUN e WALL JUMP}
\begin{oggetti}
24 & Pedana & P5\_\allowbreak{}WallRun & (5700; 0; 930) & (500; 500; 40) & MI\_\allowbreak{}AIV\_\allowbreak{}Percorso, si cammina a Z = 950 \\
25 & Text Render & Scritta & (5650; 0; 1230) & yaw 180 & ``WALL RUN: corri e salta verso il muro blu / salta di nuovo per passare all'altro muro'', World Size 32, colore (0,5; 0,8; 1) \\
26 & Cubo & Muro\_\allowbreak{}A & (6450; 350; 1100) & (800; 100; 700) & MI\_\allowbreak{}AIV\_\allowbreak{}Muro, Tag: WallRun \\
27 & Cubo & Muro\_\allowbreak{}B & (7300; -350; 1100) & (800; 100; 700) & MI\_\allowbreak{}AIV\_\allowbreak{}Muro, Tag: WallRun \\
28 & BP\_\allowbreak{}Moneta & Moneta\_\allowbreak{}MuroA & (6500; 220; 1140) & -- & Value = 1; Material = MI\_\allowbreak{}AIV\_\allowbreak{}Moneta; GlowIntensity = 4; GlowColor = (1; 0,75; 0,3) \\
29 & BP\_\allowbreak{}Moneta & Moneta\_\allowbreak{}MuroB & (7350; -220; 1140) & -- & Value = 1; Material = MI\_\allowbreak{}AIV\_\allowbreak{}Moneta; GlowIntensity = 4; GlowColor = (1; 0,75; 0,3) \\
\nota{se cadi: un pavimento di sicurezza e una pedana che ti rilancia su P5}
30 & Pedana & Sotto\_\allowbreak{}Sicurezza & (6900; 0; 480) & (2000; 1200; 40) & MI\_\allowbreak{}AIV\_\allowbreak{}PercorsoScuro, si cammina a Z = 500 \\
31 & BP\_\allowbreak{}Trampolino & Pedana\_\allowbreak{}Ritorno & (6650; 0; 510) & yaw 180 & LaunchVelocity = (350; 0; 1300); OverrideHorizontal = sì; PadMaterial = MI\_\allowbreak{}AIV\_\allowbreak{}Salto \\
32 & Text Render & Scritta & (6650; 0; 760) & yaw 180 & ``caduto? la pedana ti riporta su'', World Size 30, colore (1; 1; 1) \\
33 & Pedana & P6\_\allowbreak{}DopoMuri & (8150; 0; 930) & (800; 1000; 40) & MI\_\allowbreak{}AIV\_\allowbreak{}Percorso, si cammina a Z = 950 \\
\end{oggetti}

\subsection{Corridoio di lava (toglie vita)}
\begin{oggetti}
34 & Pedana & P7\_\allowbreak{}Lava & (9250; 0; 930) & (1400; 600; 40) & MI\_\allowbreak{}AIV\_\allowbreak{}PercorsoScuro, si cammina a Z = 950 \\
35 & BP\_\allowbreak{}ZonaDanno & Lava & (9250; 0; 1010) & -- & ZoneExtent = (400; 300; 60); DamagePerHit = 10; HitInterval = 1; Label = ``LAVA: 10 HP al secondo, attraversa in fretta''; VisualMaterial = MI\_\allowbreak{}AIV\_\allowbreak{}Lava; VisualThickness = 6 \\
36 & Point Light & Luce\_\allowbreak{}Lava & (9250; 0; 1080) & -- & 120 cd, raggio 900, colore (1; 0,35; 0,08), niente ombre \\
37 & Decal & Bruciatura\_\allowbreak{}0 & (8780; 180; 960) & rot (-90; 0; 0), Decal Size (30; 130; 130) & MI\_\allowbreak{}AIV\_\allowbreak{}Bruciatura \\
38 & Decal & Bruciatura\_\allowbreak{}1 & (8800; -200; 960) & rot (-90; 40; 0), Decal Size (30; 130; 130) & MI\_\allowbreak{}AIV\_\allowbreak{}Bruciatura \\
39 & Decal & Bruciatura\_\allowbreak{}2 & (9720; 120; 960) & rot (-90; 80; 0), Decal Size (30; 130; 130) & MI\_\allowbreak{}AIV\_\allowbreak{}Bruciatura \\
40 & Decal & Bruciatura\_\allowbreak{}3 & (9700; -170; 960) & rot (-90; 120; 0), Decal Size (30; 130; 130) & MI\_\allowbreak{}AIV\_\allowbreak{}Bruciatura \\
41 & BP\_\allowbreak{}Moneta & Moneta\_\allowbreak{}Lava\_\allowbreak{}1 & (9050; 0; 1040) & -- & Value = 1; Material = MI\_\allowbreak{}AIV\_\allowbreak{}Moneta; GlowIntensity = 4; GlowColor = (1; 0,75; 0,3) \\
42 & BP\_\allowbreak{}Moneta & Moneta\_\allowbreak{}Lava\_\allowbreak{}2 & (9450; 0; 1040) & -- & Value = 1; Material = MI\_\allowbreak{}AIV\_\allowbreak{}Moneta; GlowIntensity = 4; GlowColor = (1; 0,75; 0,3) \\
\end{oggetti}

\subsection{La sbarra rotante che spinge}
\begin{oggetti}
43 & Pedana & P8\_\allowbreak{}Sbarra & (10450; 0; 930) & (800; 800; 40) & MI\_\allowbreak{}AIV\_\allowbreak{}Percorso, si cammina a Z = 950 \\
44 & BP\_\allowbreak{}BarraRotante & Sbarra\_\allowbreak{}Rotante & (10450; 0; 950) & -- & DegreesPerSecond = 75; BarSize = (650; 30; 40); BarHeight = 60; PushStrength = 900; KillOnTouch = no \\
45 & BP\_\allowbreak{}Moneta & Moneta\_\allowbreak{}Sbarra & (10450; 300; 1040) & -- & Value = 1; Material = MI\_\allowbreak{}AIV\_\allowbreak{}Moneta; GlowIntensity = 4; GlowColor = (1; 0,75; 0,3) \\
\end{oggetti}

\subsection{Ghiaccio e negozio}
\begin{oggetti}
46 & Pedana & P9\_\allowbreak{}Piazza & (11400; 0; 930) & (900; 1200; 40) & MI\_\allowbreak{}AIV\_\allowbreak{}PercorsoScuro, si cammina a Z = 950 \\
47 & Cubo & Ponte\_\allowbreak{}Ghiaccio & (12350; -250; 930) & (1000; 300; 40) & MI\_\allowbreak{}AIV\_\allowbreak{}Ghiaccio \\
48 & BP\_\allowbreak{}ZonaModificatore & Zona\_\allowbreak{}Ghiaccio & (12075; -250; 1050) & -- & ZoneExtent = (775; 150; 100); GroundFrictionMultiplier = 0,08; BrakingMultiplier = 0,05; DurationAfterExit = 0,5; Label = ``GHIACCIO''; ShowVisual = no \\
49 & BP\_\allowbreak{}Ascensore & Ponte\_\allowbreak{}Pagato & (12350; 250; 535) & -- & EndOffset = (0; 0; 400); PlatformSize = (1000; 300; 30); TransitionTime = 2; Material (mesh) = MI\_\allowbreak{}AIV\_\allowbreak{}Moneta \\
50 & BP\_\allowbreak{}Negozio & Negozio & (11600; 400; 1050) & -- & ZoneExtent = (120; 120; 100); Cost = 5; OnlyOnce = sì; Label = ``5 MONETE: ponte sicuro''; VisualMaterial = MI\_\allowbreak{}AIV\_\allowbreak{}Moneta; Targets = [Ponte\_\allowbreak{}Pagato] \\
51 & Text Render & Scritta & (11150; 0; 1250) & yaw 180 & ``ghiaccio gratis... o ponte sicuro a 5 monete'', World Size 32, colore (1; 1; 1) \\
\end{oggetti}

\subsection{JETPACK: il salto lunghissimo verso una piattaforma piccola}
\begin{oggetti}
52 & Pedana & P10\_\allowbreak{}Jetpack & (13250; 0; 930) & (800; 1000; 40) & MI\_\allowbreak{}AIV\_\allowbreak{}Salto, si cammina a Z = 950 \\
53 & BP\_\allowbreak{}PickupAbilita & Oggetto\_\allowbreak{}Jetpack & (13100; 0; 1060) & -- & Ability = Jetpack; Material = MI\_\allowbreak{}AIV\_\allowbreak{}Abilita; GlowIntensity = 20; GlowColor = (0,3; 0,7; 1); PickupRadius = 80; Label = ``JETPACK''; Message = ``Jetpack! In aria TIENI PREMUTO salto: voli finche' c'e' carburante'' \\
54 & Spot Light & Luce\_\allowbreak{}Oggetto\_\allowbreak{}Jetpack & (13100; 0; 1650) & pitch -90 yaw 0 & 300 cd, raggio 1000, colore (0,5; 0,75; 1), coni 12/22° \\
55 & Text Render & Scritta & (13300; 0; 1250) & yaw 180 & ``JETPACK: salta e TIENI PREMUTO salto / pilota fino alla piattaforma piccola'', World Size 32, colore (1; 0,65; 0,3) \\
56 & Pedana & P11\_\allowbreak{}Atterraggio & (14950; 0; 1080) & (300; 300; 40) & MI\_\allowbreak{}AIV\_\allowbreak{}Salto, si cammina a Z = 1100 \\
57 & BP\_\allowbreak{}Moneta & Moneta\_\allowbreak{}Atterraggio & (14950; 0; 1190) & -- & Value = 2; Material = MI\_\allowbreak{}AIV\_\allowbreak{}Moneta; GlowIntensity = 4; GlowColor = (1; 0,75; 0,3) \\
\end{oggetti}

\subsection{Piattaforme che si restringono mentre ti avvicini}
\begin{oggetti}
58 & BP\_\allowbreak{}CresceConDistanza & Si\_\allowbreak{}Restringe\_\allowbreak{}1 & (15700; 0; 1085) & -- & Material = MI\_\allowbreak{}AIV\_\allowbreak{}Salto; BaseSize = (400; 400; 30); MinScale = 0,45; MaxScale = 1; NearDistance = 150; FarDistance = 900; ScaleHeight = no; Label = ``SI RESTRINGE QUANDO TI AVVICINI'' \\
59 & BP\_\allowbreak{}CresceConDistanza & Si\_\allowbreak{}Restringe\_\allowbreak{}2 & (16400; 250; 1085) & -- & Material = MI\_\allowbreak{}AIV\_\allowbreak{}Salto; BaseSize = (400; 400; 30); MinScale = 0,45; MaxScale = 1; NearDistance = 150; FarDistance = 900; ScaleHeight = no \\
60 & BP\_\allowbreak{}CresceConDistanza & Si\_\allowbreak{}Restringe\_\allowbreak{}3 & (17100; -150; 1085) & -- & Material = MI\_\allowbreak{}AIV\_\allowbreak{}Salto; BaseSize = (400; 400; 30); MinScale = 0,45; MaxScale = 1; NearDistance = 150; FarDistance = 900; ScaleHeight = no \\
\end{oggetti}

\subsection{Torretta con i proiettili che fanno RIPARTIRE il livello}
\begin{oggetti}
61 & Pedana & P12\_\allowbreak{}TiroMortale & (18300; 0; 1080) & (1400; 700; 40) & MI\_\allowbreak{}AIV\_\allowbreak{}Percorso, si cammina a Z = 1100 \\
62 & Cubo & Riparo\_\allowbreak{}Mortale\_\allowbreak{}1 & (17950; -120; 1200) & (80; 300; 200) & MI\_\allowbreak{}AIV\_\allowbreak{}LegnoScuro \\
63 & Cubo & Riparo\_\allowbreak{}Mortale\_\allowbreak{}2 & (18400; 150; 1200) & (80; 300; 200) & MI\_\allowbreak{}AIV\_\allowbreak{}LegnoScuro \\
64 & Cubo & Riparo\_\allowbreak{}Mortale\_\allowbreak{}3 & (18800; -120; 1200) & (80; 300; 200) & MI\_\allowbreak{}AIV\_\allowbreak{}LegnoScuro \\
65 & Pedana & Torretta\_\allowbreak{}Mortale\_\allowbreak{}Base & (18300; 900; 750) & (250; 250; 900) & MI\_\allowbreak{}AIV\_\allowbreak{}Pericolo, si cammina a Z = 1200 \\
66 & BP\_\allowbreak{}Torretta & Torretta\_\allowbreak{}Mortale & (18300; 900; 1200) & yaw -90 & Range = 1600; FireInterval = 1,8; TurnSpeed = 2,5; ProjectileSpeed = 1100; RestartLevelOnHit = sì; BodyMaterial = MI\_\allowbreak{}AIV\_\allowbreak{}Pericolo; ProjectileMaterial = MI\_\allowbreak{}AIV\_\allowbreak{}Proiettile \\
67 & Text Render & Scritta & (17700; 0; 1400) & yaw 180 & ``PROIETTILI MORTALI: se ti colpiscono il livello riparte / vai da un riparo all'altro'', World Size 32, colore (1; 0,35; 0,3) \\
68 & Pedana & P13\_\allowbreak{}Arrivo & (19500; 0; 1080) & (800; 800; 40) & MI\_\allowbreak{}AIV\_\allowbreak{}Traguardo, si cammina a Z = 1100 \\
69 & BP\_\allowbreak{}Traguardo & Traguardo & (19550; 0; 1200) & -- & ZoneExtent = (150; 150; 100); NextLevel = L\_\allowbreak{}GD13\_\allowbreak{}Livello3\_\allowbreak{}BP; Message = ``Livello 2 completato! Ora di notte, a tempo''; Delay = 1,2; ResetDeathCount = no; Label = ``TRAGUARDO''; VisualMaterial = MI\_\allowbreak{}AIV\_\allowbreak{}Traguardo \\
70 & Point Light & Luce\_\allowbreak{}Traguardo & (19550; 0; 1360) & -- & 30 cd, raggio 700, colore (0,4; 1; 0,5), niente ombre \\
\end{oggetti}



\textbf{Collaudo del livello 2.} (1) Il trampolino porta su P2. (2) La leva (\tasto{E}) fa salire l'ascensore fino a P5. (3) Si corre sul Muro\_A, si salta sul Muro\_B, si arriva a P6; cadendo, la pedana di ritorno rilancia su P5. (4) Nella lava la vita scende di 10 al secondo. (5) Con 5 monete il negozio alza il ponte. (6) Con il jetpack (salto tenuto premuto) si arriva a P11 da P10. (7) Un proiettile della torretta mortale fa ripartire il livello.

% =====================================================================
\section{Livello 3 -- ``Di notte, a tempo''}

Notte con la luna e la nebbia, 150 secondi. Piattaforma laterale, portale, gravità ridotta, sbarra che uccide (si salta), wall run, torretta che anticipa il movimento, il tunnel buio da fare in prima persona con la torcia, la piattaforma che cambia forma ogni secondo e la combo finale: wall run, jetpack, casse e doppio salto. Il traguardo riporta al menu e azzera le morti.

Livello: \var{L\_GD13\_Livello3\_BP}. Traguardo $\rightarrow$ \var{L\_GD13\_Menu}.

%% generato da strumenti/estrai_livelli.py: non modificare a mano
\subsection{Luce, atmosfera e partenza}
\begin{itemize}
\item \nodo{Directional Light} (la luna): Mobility \emph{Movable}, Rotation pitch -35, yaw -120; Intensity 0,4 lux; \emph{Use Temperature} sì, Temperature 9000 K; \emph{Atmosphere Sun Light} sì; Source Angle 0,5.
\item \nodo{Sky Atmosphere} (valori di default).
\item \nodo{Sky Light}: Movable, \emph{Real Time Capture} sì, Intensity 0,6.
\item \nodo{Exponential Height Fog} a Z $-200$: Fog Density 0,03, \emph{Volumetric Fog} sì.
\item \nodo{Post Process Volume}: \emph{Infinite Extent (Unbound)} sì; spuntate e impostate Auto Exposure Bias 1, Auto Exposure MinBrightness 0,4, Auto Exposure MaxBrightness 1,5, Bloom Intensity 0,9, Vignette Intensity 0,45.
\item \nodo{Player Start} in (-150; 0; 120), rotazione 0.
\end{itemize}

\subsection{Partenza e primo tratto}
\begin{oggetti}
1 & BP\_\allowbreak{}ZonaCaduta & Zona\_\allowbreak{}Caduta & (9000; 0; -600) & -- & ZoneExtent = (12000; 4000; 150); ShowVisual = no \\
\nota{il timer: se scade, il livello riparte (lo mostra anche l'HUD)}
2 & BP\_\allowbreak{}TimerLivello & Timer & (300; 0; 450) & yaw 180 & TimeLimit = 150; WarningThreshold = 20; ShowOnScreen = no; ShowInWorld = sì \\
3 & Text Render & Scritta & (300; 0; 300) & yaw 180 & ``LIVELLO 3: 150 secondi'', World Size 45, colore (1; 1; 1) \\
4 & Pedana & P0\_\allowbreak{}Partenza & (0; 0; -20) & (600; 600; 40) & MI\_\allowbreak{}AIV\_\allowbreak{}Percorso, si cammina a Z = 0 \\
5 & Cubo & Braciere\_\allowbreak{}0\_\allowbreak{}Base & (200; 230; 45) & (50; 50; 90) & MI\_\allowbreak{}AIV\_\allowbreak{}Ferro \\
6 & Static Mesh & Braciere\_\allowbreak{}0\_\allowbreak{}Braci & (200; 230; 95) & scala (0,45; 0,45; 0,22) & Sphere; MI\_\allowbreak{}AIV\_\allowbreak{}Braci, Cast Shadow no, Collision: No Collision \\
7 & BP\_\allowbreak{}LuceTremolante & Braciere\_\allowbreak{}0\_\allowbreak{}Luce & (200; 230; 150) & -- & BaseIntensity = 80; FlickerAmount = 0,3; FlickerSpeed = 7; Temperature = 1800; AttenuationRadius = 1400; SourceRadius = 10; PositionJitter = 5; GlowActors = [Braciere\_\allowbreak{}0\_\allowbreak{}Braci] \\
\nota{piattaforma laterale}
8 & BP\_\allowbreak{}PiattaformaMobile & Mobile\_\allowbreak{}Laterale & (550; -350; -15) & -- & EndOffset = (0; 700; 0); TravelTime = 2; PauseTime = 0,4; StartDelay = 0; PlatformSize = (250; 250; 30); PlatformMaterial = MI\_\allowbreak{}AIV\_\allowbreak{}Salto; MoveSound = SfxPlatformMoving\_\allowbreak{}Cue \\
9 & Pedana & P1\_\allowbreak{}Portale & (1000; 0; -20) & (300; 300; 40) & MI\_\allowbreak{}AIV\_\allowbreak{}Salto, si cammina a Z = 0 \\
\nota{portale: teletrasporto su un'isola alta}
10 & BP\_\allowbreak{}Teletrasporto & Portale & (1000; 0; 100) & -- & ZoneExtent = (100; 100; 100); Destination = (1500; 0; 600); StopVelocity = sì; Label = ``PORTALE''; VisualMaterial = MI\_\allowbreak{}AIV\_\allowbreak{}Abilita \\
11 & Point Light & Luce\_\allowbreak{}Portale & (1000; 0; 180) & -- & 60 cd, raggio 600, colore (0,35; 0,6; 1), niente ombre \\
\nota{gravità ridotta per il salto lungo}
12 & Pedana & P2\_\allowbreak{}Luna & (2500; 0; 580) & (500; 500; 40) & MI\_\allowbreak{}AIV\_\allowbreak{}Percorso, si cammina a Z = 600 \\
13 & Cubo & Braciere\_\allowbreak{}2\_\allowbreak{}Base & (2350; -200; 645) & (50; 50; 90) & MI\_\allowbreak{}AIV\_\allowbreak{}Ferro \\
14 & Static Mesh & Braciere\_\allowbreak{}2\_\allowbreak{}Braci & (2350; -200; 695) & scala (0,45; 0,45; 0,22) & Sphere; MI\_\allowbreak{}AIV\_\allowbreak{}Braci, Cast Shadow no, Collision: No Collision \\
15 & BP\_\allowbreak{}LuceTremolante & Braciere\_\allowbreak{}2\_\allowbreak{}Luce & (2350; -200; 750) & -- & BaseIntensity = 80; FlickerAmount = 0,3; FlickerSpeed = 7; Temperature = 1800; AttenuationRadius = 1400; SourceRadius = 10; PositionJitter = 5; GlowActors = [Braciere\_\allowbreak{}2\_\allowbreak{}Braci] \\
16 & BP\_\allowbreak{}ZonaModificatore & Zona\_\allowbreak{}GravitaRidotta & (2600; 0; 700) & -- & ZoneExtent = (150; 200; 100); WalkSpeedMultiplier = 1,1; JumpMultiplier = 1,2; GravityMultiplier = 0,45; DurationAfterExit = 3; Label = ``GRAVITA RIDOTTA''; VisualMaterial = MI\_\allowbreak{}AIV\_\allowbreak{}Salto \\
17 & Pedana & P3 & (3700; 0; 680) & (400; 400; 40) & MI\_\allowbreak{}AIV\_\allowbreak{}PercorsoScuro, si cammina a Z = 700 \\
18 & BP\_\allowbreak{}Moneta & Moneta\_\allowbreak{}P3 & (3700; 0; 790) & -- & Value = 1; Material = MI\_\allowbreak{}AIV\_\allowbreak{}Moneta; GlowIntensity = 4; GlowColor = (1; 0,75; 0,3) \\
\nota{sbarra che uccide: va saltata}
19 & Pedana & P4\_\allowbreak{}SbarraMortale & (4600; 0; 680) & (800; 800; 40) & MI\_\allowbreak{}AIV\_\allowbreak{}Percorso, si cammina a Z = 700 \\
20 & BP\_\allowbreak{}BarraRotante & Sbarra\_\allowbreak{}Mortale & (4600; 0; 700) & -- & DegreesPerSecond = 60; BarSize = (600; 30; 30); BarHeight = 50; KillOnTouch = sì \\
21 & Cubo & Braciere\_\allowbreak{}4\_\allowbreak{}Base & (4950; 350; 745) & (50; 50; 90) & MI\_\allowbreak{}AIV\_\allowbreak{}Ferro \\
22 & Static Mesh & Braciere\_\allowbreak{}4\_\allowbreak{}Braci & (4950; 350; 795) & scala (0,45; 0,45; 0,22) & Sphere; MI\_\allowbreak{}AIV\_\allowbreak{}Braci, Cast Shadow no, Collision: No Collision \\
23 & BP\_\allowbreak{}LuceTremolante & Braciere\_\allowbreak{}4\_\allowbreak{}Luce & (4950; 350; 850) & -- & BaseIntensity = 80; FlickerAmount = 0,3; FlickerSpeed = 7; Temperature = 1800; AttenuationRadius = 1400; SourceRadius = 10; PositionJitter = 5; GlowActors = [Braciere\_\allowbreak{}4\_\allowbreak{}Braci] \\
\nota{wall run veloce (oppure un doppio salto ben fatto)}
24 & Cubo & Muro\_\allowbreak{}Notte & (5500; 350; 850) & (900; 100; 700) & MI\_\allowbreak{}AIV\_\allowbreak{}Muro, Tag: WallRun \\
25 & BP\_\allowbreak{}Moneta & Moneta\_\allowbreak{}Muro & (5500; 220; 890) & -- & Value = 1; Material = MI\_\allowbreak{}AIV\_\allowbreak{}Moneta; GlowIntensity = 4; GlowColor = (1; 0,75; 0,3) \\
26 & Pedana & P5 & (6400; 0; 680) & (800; 800; 40) & MI\_\allowbreak{}AIV\_\allowbreak{}Percorso, si cammina a Z = 700 \\
27 & Cubo & Braciere\_\allowbreak{}5\_\allowbreak{}Base & (6700; -300; 745) & (50; 50; 90) & MI\_\allowbreak{}AIV\_\allowbreak{}Ferro \\
28 & Static Mesh & Braciere\_\allowbreak{}5\_\allowbreak{}Braci & (6700; -300; 795) & scala (0,45; 0,45; 0,22) & Sphere; MI\_\allowbreak{}AIV\_\allowbreak{}Braci, Cast Shadow no, Collision: No Collision \\
29 & BP\_\allowbreak{}LuceTremolante & Braciere\_\allowbreak{}5\_\allowbreak{}Luce & (6700; -300; 850) & -- & BaseIntensity = 80; FlickerAmount = 0,3; FlickerSpeed = 7; Temperature = 1800; AttenuationRadius = 1400; SourceRadius = 10; PositionJitter = 5; GlowActors = [Braciere\_\allowbreak{}5\_\allowbreak{}Braci] \\
30 & BP\_\allowbreak{}PickupSalute & Cura\_\allowbreak{}P5 & (6300; 250; 790) & -- & HealAmount = 40; Material = MI\_\allowbreak{}AIV\_\allowbreak{}Vita; GlowColor = (0,3; 1; 0,4); Label = ``+VITA'' \\
\nota{salti da doppio salto}
31 & Pedana & P6 & (7400; 0; 880) & (300; 300; 40) & MI\_\allowbreak{}AIV\_\allowbreak{}Salto, si cammina a Z = 900 \\
32 & Pedana & P7 & (8200; 0; 980) & (300; 300; 40) & MI\_\allowbreak{}AIV\_\allowbreak{}Salto, si cammina a Z = 1000 \\
33 & BP\_\allowbreak{}Moneta & Moneta\_\allowbreak{}P6 & (7400; 0; 990) & -- & Value = 1; Material = MI\_\allowbreak{}AIV\_\allowbreak{}Moneta; GlowIntensity = 4; GlowColor = (1; 0,75; 0,3) \\
34 & BP\_\allowbreak{}Moneta & Moneta\_\allowbreak{}P7 & (8200; 0; 1090) & -- & Value = 1; Material = MI\_\allowbreak{}AIV\_\allowbreak{}Moneta; GlowIntensity = 4; GlowColor = (1; 0,75; 0,3) \\
\nota{torretta che anticipa il movimento}
35 & Pedana & P8 & (9000; 0; 980) & (600; 600; 40) & MI\_\allowbreak{}AIV\_\allowbreak{}Percorso, si cammina a Z = 1000 \\
36 & Cubo & Riparo\_\allowbreak{}P8 & (9000; 120; 1100) & (80; 260; 200) & MI\_\allowbreak{}AIV\_\allowbreak{}LegnoScuro \\
37 & Pedana & Torretta\_\allowbreak{}Base & (8800; 700; 700) & (250; 250; 800) & MI\_\allowbreak{}AIV\_\allowbreak{}Pericolo, si cammina a Z = 1100 \\
38 & BP\_\allowbreak{}Torretta & Torretta\_\allowbreak{}Furba & (8800; 700; 1100) & yaw -90 & Range = 1600; FireInterval = 1,8; ProjectileSpeed = 1600; ProjectileDamage = 20; LeadTarget = sì; BodyMaterial = MI\_\allowbreak{}AIV\_\allowbreak{}Torretta; ProjectileMaterial = MI\_\allowbreak{}AIV\_\allowbreak{}Proiettile \\
39 & Cubo & Braciere\_\allowbreak{}8\_\allowbreak{}Base & (9200; -230; 1045) & (50; 50; 90) & MI\_\allowbreak{}AIV\_\allowbreak{}Ferro \\
40 & Static Mesh & Braciere\_\allowbreak{}8\_\allowbreak{}Braci & (9200; -230; 1095) & scala (0,45; 0,45; 0,22) & Sphere; MI\_\allowbreak{}AIV\_\allowbreak{}Braci, Cast Shadow no, Collision: No Collision \\
41 & BP\_\allowbreak{}LuceTremolante & Braciere\_\allowbreak{}8\_\allowbreak{}Luce & (9200; -230; 1150) & -- & BaseIntensity = 80; FlickerAmount = 0,3; FlickerSpeed = 7; Temperature = 1800; AttenuationRadius = 1400; SourceRadius = 10; PositionJitter = 5; GlowActors = [Braciere\_\allowbreak{}8\_\allowbreak{}Braci] \\
\nota{ascensore verso il tunnel}
42 & BP\_\allowbreak{}PiattaformaMobile & Ascensore & (9550; 0; 985) & -- & EndOffset = (0; 0; 500); TravelTime = 2; PauseTime = 1; StartDelay = 0; PlatformSize = (250; 250; 30); PlatformMaterial = MI\_\allowbreak{}AIV\_\allowbreak{}Salto; MoveSound = SfxPlatformMoving\_\allowbreak{}Cue \\
\end{oggetti}

\subsection{IL TUNNEL BUIO: si passa solo con la torcia (prima persona)}
\begin{oggetti}
43 & Pedana & P9\_\allowbreak{}Tunnel\_\allowbreak{}Ingresso & (10000; 0; 1480) & (500; 500; 40) & MI\_\allowbreak{}AIV\_\allowbreak{}PercorsoScuro, si cammina a Z = 1500 \\
44 & Cubo & Braciere\_\allowbreak{}9\_\allowbreak{}Base & (9850; 200; 1545) & (50; 50; 90) & MI\_\allowbreak{}AIV\_\allowbreak{}Ferro \\
45 & Static Mesh & Braciere\_\allowbreak{}9\_\allowbreak{}Braci & (9850; 200; 1595) & scala (0,45; 0,45; 0,22) & Sphere; MI\_\allowbreak{}AIV\_\allowbreak{}Braci, Cast Shadow no, Collision: No Collision \\
46 & BP\_\allowbreak{}LuceTremolante & Braciere\_\allowbreak{}9\_\allowbreak{}Luce & (9850; 200; 1650) & -- & BaseIntensity = 100; FlickerAmount = 0,3; FlickerSpeed = 7; Temperature = 1800; AttenuationRadius = 1400; SourceRadius = 10; PositionJitter = 5; GlowActors = [Braciere\_\allowbreak{}9\_\allowbreak{}Braci] \\
47 & BP\_\allowbreak{}PickupBatteria & Batteria\_\allowbreak{}Ingresso & (10050; -150; 1590) & -- & BatterySeconds = 15; Material = MI\_\allowbreak{}AIV\_\allowbreak{}Moneta; Label = ``BATTERIA'' \\
48 & Text Render & Scritta & (10000; 0; 1830) & yaw 180 & ``TUNNEL BUIO / rotellina del mouse AVANTI: prima persona e torcia / (la batteria dura poco: le batterie la ricaricano)'', World Size 28, colore (1; 0,95; 0,5) \\
49 & Cubo & Tunnel\_\allowbreak{}Tetto & (11400; 0; 1940) & (2200; 840; 40) & MI\_\allowbreak{}AIV\_\allowbreak{}PercorsoScuro \\
50 & Cubo & Tunnel\_\allowbreak{}Parete\_\allowbreak{}S & (11400; -420; 1570) & (2200; 40; 780) & MI\_\allowbreak{}AIV\_\allowbreak{}PercorsoScuro \\
51 & Cubo & Tunnel\_\allowbreak{}Parete\_\allowbreak{}N & (11400; 420; 1570) & (2200; 40; 780) & MI\_\allowbreak{}AIV\_\allowbreak{}PercorsoScuro \\
52 & Pedana & Tunnel\_\allowbreak{}P1 & (10500; 120; 1480) & (260; 260; 40) & MI\_\allowbreak{}AIV\_\allowbreak{}PercorsoScuro, si cammina a Z = 1500 \\
53 & Pedana & Tunnel\_\allowbreak{}P2 & (10950; -120; 1480) & (260; 260; 40) & MI\_\allowbreak{}AIV\_\allowbreak{}PercorsoScuro, si cammina a Z = 1500 \\
54 & Pedana & Tunnel\_\allowbreak{}P3 & (11400; 120; 1480) & (260; 260; 40) & MI\_\allowbreak{}AIV\_\allowbreak{}PercorsoScuro, si cammina a Z = 1500 \\
55 & Pedana & Tunnel\_\allowbreak{}P4 & (11850; -120; 1480) & (260; 260; 40) & MI\_\allowbreak{}AIV\_\allowbreak{}PercorsoScuro, si cammina a Z = 1500 \\
56 & Pedana & Tunnel\_\allowbreak{}P5 & (12300; 120; 1480) & (260; 260; 40) & MI\_\allowbreak{}AIV\_\allowbreak{}PercorsoScuro, si cammina a Z = 1500 \\
57 & BP\_\allowbreak{}PickupBatteria & Batteria\_\allowbreak{}Tunnel & (11400; 120; 1590) & -- & BatterySeconds = 15; Material = MI\_\allowbreak{}AIV\_\allowbreak{}Moneta; Label = ``BATTERIA'' \\
58 & BP\_\allowbreak{}Moneta & Moneta\_\allowbreak{}Tunnel & (11850; -120; 1590) & -- & Value = 3; Material = MI\_\allowbreak{}AIV\_\allowbreak{}Moneta; GlowIntensity = 4; GlowColor = (1; 0,75; 0,3) \\
\end{oggetti}

\subsection{La piattaforma che cambia modello ogni secondo}
\begin{oggetti}
59 & Pedana & P10\_\allowbreak{}DopoTunnel & (12900; 0; 1480) & (500; 500; 40) & MI\_\allowbreak{}AIV\_\allowbreak{}Percorso, si cammina a Z = 1500 \\
60 & Cubo & Braciere\_\allowbreak{}10\_\allowbreak{}Base & (12800; 200; 1545) & (50; 50; 90) & MI\_\allowbreak{}AIV\_\allowbreak{}Ferro \\
61 & Static Mesh & Braciere\_\allowbreak{}10\_\allowbreak{}Braci & (12800; 200; 1595) & scala (0,45; 0,45; 0,22) & Sphere; MI\_\allowbreak{}AIV\_\allowbreak{}Braci, Cast Shadow no, Collision: No Collision \\
62 & BP\_\allowbreak{}LuceTremolante & Braciere\_\allowbreak{}10\_\allowbreak{}Luce & (12800; 200; 1650) & -- & BaseIntensity = 80; FlickerAmount = 0,3; FlickerSpeed = 7; Temperature = 1800; AttenuationRadius = 1400; SourceRadius = 10; PositionJitter = 5; GlowActors = [Braciere\_\allowbreak{}10\_\allowbreak{}Braci] \\
63 & BP\_\allowbreak{}CambiaModello & Piattaforma\_\allowbreak{}CambiaForma & (13550; 0; 1400) & -- & Interval = 1; Size = (250; 250; 200); Collide = sì; Materials = [MI\_\allowbreak{}AIV\_\allowbreak{}Salto, MI\_\allowbreak{}AIV\_\allowbreak{}Abilita, MI\_\allowbreak{}AIV\_\allowbreak{}Pericolo, MI\_\allowbreak{}AIV\_\allowbreak{}Moneta]; Label = ``CAMBIA FORMA OGNI SECONDO: salta quando e' un cubo'' \\
64 & Pedana & P11\_\allowbreak{}Combo & (14200; 0; 1480) & (500; 500; 40) & MI\_\allowbreak{}AIV\_\allowbreak{}Percorso, si cammina a Z = 1500 \\
\end{oggetti}

\subsection{COMBO FINALE: wall run $\rightarrow$ jetpack $\rightarrow$ casse + doppio salto}
\begin{oggetti}
65 & Cubo & Muro\_\allowbreak{}Combo & (15000; 350; 1650) & (900; 100; 700) & MI\_\allowbreak{}AIV\_\allowbreak{}Muro, Tag: WallRun \\
66 & Text Render & Scritta & (14200; 0; 1800) & yaw 180 & ``COMBO: wall run, poi jetpack fino alla piattaforma, / poi casse e doppio salto'', World Size 32, colore (0,5; 0,8; 1) \\
67 & Pedana & P12\_\allowbreak{}Casse & (16700; 0; 1580) & (600; 600; 40) & MI\_\allowbreak{}AIV\_\allowbreak{}Salto, si cammina a Z = 1600 \\
68 & Cubo & Braciere\_\allowbreak{}12\_\allowbreak{}Base & (16550; -230; 1645) & (50; 50; 90) & MI\_\allowbreak{}AIV\_\allowbreak{}Ferro \\
69 & Static Mesh & Braciere\_\allowbreak{}12\_\allowbreak{}Braci & (16550; -230; 1695) & scala (0,45; 0,45; 0,22) & Sphere; MI\_\allowbreak{}AIV\_\allowbreak{}Braci, Cast Shadow no, Collision: No Collision \\
70 & BP\_\allowbreak{}LuceTremolante & Braciere\_\allowbreak{}12\_\allowbreak{}Luce & (16550; -230; 1750) & -- & BaseIntensity = 80; FlickerAmount = 0,3; FlickerSpeed = 7; Temperature = 1800; AttenuationRadius = 1400; SourceRadius = 10; PositionJitter = 5; GlowActors = [Braciere\_\allowbreak{}12\_\allowbreak{}Braci] \\
71 & Pedana & P13\_\allowbreak{}Arrivo & (17250; 0; 1610) & (500; 600; 900) & MI\_\allowbreak{}AIV\_\allowbreak{}Traguardo, si cammina a Z = 2060 \\
72 & Cubo & Braciere\_\allowbreak{}13\_\allowbreak{}Base & (17350; 220; 2105) & (50; 50; 90) & MI\_\allowbreak{}AIV\_\allowbreak{}Ferro \\
73 & Static Mesh & Braciere\_\allowbreak{}13\_\allowbreak{}Braci & (17350; 220; 2155) & scala (0,45; 0,45; 0,22) & Sphere; MI\_\allowbreak{}AIV\_\allowbreak{}Braci, Cast Shadow no, Collision: No Collision \\
74 & BP\_\allowbreak{}LuceTremolante & Braciere\_\allowbreak{}13\_\allowbreak{}Luce & (17350; 220; 2210) & -- & BaseIntensity = 100; FlickerAmount = 0,3; FlickerSpeed = 7; Temperature = 1800; AttenuationRadius = 1400; SourceRadius = 10; PositionJitter = 5; GlowActors = [Braciere\_\allowbreak{}13\_\allowbreak{}Braci] \\
75 & BP\_\allowbreak{}Traguardo & Traguardo & (17300; 0; 2160) & -- & ZoneExtent = (150; 150; 100); NextLevel = L\_\allowbreak{}GD13\_\allowbreak{}Menu; Message = ``Hai finito tutti e tre i livelli! Si torna al menu''; Delay = 1,2; ResetDeathCount = sì; Label = ``TRAGUARDO''; VisualMaterial = MI\_\allowbreak{}AIV\_\allowbreak{}Traguardo \\
76 & Point Light & Luce\_\allowbreak{}Traguardo & (17300; 0; 2320) & -- & 30 cd, raggio 700, colore (0,4; 1; 0,5), niente ombre \\
\end{oggetti}



\textbf{Collaudo del livello 3.} (1) Il timer parte e, a zero, il livello riparte. (2) Il portale porta su P2\_Luna; nella zona arancione si salta molto più lontano. (3) Toccare la sbarra mortale fa ripartire il livello. (4) Nel tunnel, con la rotella avanti si entra in prima persona e si accende la torcia; le batterie la ricaricano. (5) La piattaforma che cambia forma si attraversa quando è un cubo. (6) Combo: dal muro blu si salta, si tiene premuto per il jetpack fino a P12, poi casse e doppio salto fino al traguardo.

% =====================================================================
\section{Quando avete finito}

\begin{itemize}
  \item Dal menu scegliete \textbf{Blueprints}: deve partire \var{L\_GD13\_Livello1\_BP}, e i traguardi devono portare a \var{\_2\_BP}, \var{\_3\_BP} e infine al menu.
  \item Giocate le due versioni una dopo l'altra: stesse distanze, stessi tempi. Se un salto riesce in una e non nell'altra, il colpevole è quasi sempre il personaggio (Jump Z Velocity, Air Control, gravità), non il livello.
  \item Se avete usato lo script dello scheletro, la cartella \texttt{DaFare} dell'Outliner deve essere vuota.
  \item \nodo{File} $\rightarrow$ \nodo{Save All}. Con \emph{One File Per Actor} (attivo nei livelli nuovi di UE5) ogni actor è un file in \texttt{Content/\_\_ExternalActors\_\_}: nel commit vanno anche quelli.
\end{itemize}

\ninja{Le tabelle di questa guida sono generate dallo script che costruisce i livelli C++ (\texttt{strumenti/estrai\_livelli.py} legge \texttt{gd13\_platformer.py} e gli header C++). Se cambiate un livello C++ nello script, rilanciate \texttt{python3 Docs/Blueprint\_Livelli/strumenti/estrai\_livelli.py} e ricompilate questa guida: le due versioni restano allineate.}

\end{document}
